% Use lualatex
\documentclass{llncs}
  
\usepackage[T1]{fontenc}
\usepackage[utf8]{inputenc}
\usepackage[russian,english]{babel}
\usepackage{fontspec}
\usepackage{lmodern}
\usepackage{mathtools}
%\usepackage{fullpage}
\usepackage{graphicx}
\usepackage{xspace}
\usepackage{tabularx}
\usepackage[lf]{ebgaramond}
\usepackage{biolinum}
\usepackage[cmintegrals,cmbraces]{newtxmath}
\usepackage{ebgaramond-maths}
\usepackage{multicol}
\usepackage{sectsty}
\usepackage[noend]{algpseudocode}
\usepackage{algorithm}
\usepackage{algorithmicx}
\usepackage[bottom]{footmisc}
\usepackage{caption}
%\captionsetup{font=small}
\captionsetup{labelfont={sf,bf}}
\usepackage{ragged2e}
\usepackage{xcolor}
\usepackage{pgfplots}
\pgfplotsset{compat=1.18}

% ebgaramond-maths, prior to version 1.3, lacks the "mu" in math. This
% hack uses the "mu" from text italics.
\DeclareFontFamily{U}{EBSUB}{}
\DeclareFontShape{U}{EBSUB}{m}{it}{
 <-> EBGaramond-Italic-tlf-lgr
}{}
\DeclareFontFamily{U}{EBSUB}{}
\DeclareSymbolFont{EBSUB}{U}{EBSUB}{m}{it}
\DeclareMathSymbol{\mu}{\mathord}{EBSUB}{`m}
\DeclareMathSymbol{\varDelta}{\mathord}{EBSUB}{`D}
\DeclareMathSymbol{\varOmega}{\mathord}{EBSUB}{`W}

\definecolor{codebg}{rgb}{.97,.97,.99}
\definecolor{gray}{rgb}{.4,.4,.4}
\definecolor{webgreen}{rgb}{0,.5,0}
\definecolor{Maroon}{cmyk}{0, 0.87, 0.68, 0.32}
\definecolor{RoyalBlue}{cmyk}{1, 0.50, 0, 0}
\definecolor{blood}{rgb}{.6,0,0}

\usepackage{listings}
\AtBeginDocument{%\counterwithin{lstlisting}{section}%
    \renewcommand\thelstlisting{\arabic{lstlisting}}}
\renewcommand{\lstlistingname}{Listing}
\DeclareCaptionFormat{listing}{#1#2#3}
\captionsetup[lstlisting]{format=listing,singlelinecheck=false}

\lstdefinestyle{zzpython}{%
    frame=ltb,
    framerule=0pt,
    aboveskip=3mm,
    belowskip=3mm,
    framextopmargin=.5mm,
    framexbottommargin=.5mm,
    framexleftmargin=1mm,
    framexrightmargin=1.5mm,
    showstringspaces=false,
    language=Python,
    columns=fullflexible,
    keepspaces=true,
    %basicstyle={\small\ttfamily},
    basicstyle={\fontsize{8pt}{9pt}\ttfamily},
    numbers=left,
    numberstyle={\fontsize{8pt}{9pt}\ttfamily\color{gray}},
    keywordstyle=\color{webgreen},
    commentstyle=\color{gray},
    stringstyle=\color{Maroon},
    identifierstyle=\color{RoyalBlue},
    breaklines=true,
    breakatwhitespace=true,
    tabsize=3,
    backgroundcolor=\color{codebg},
    captionpos=t,
}

\lstdefinestyle{zzc}{%
    frame=ltb,
    framerule=0pt,
    aboveskip=3mm,
    belowskip=3mm,
    framextopmargin=.5mm,
    framexbottommargin=.5mm,
    framexleftmargin=1mm,
    framexrightmargin=1.5mm,
    showstringspaces=false,
    language=C,
    columns=fullflexible,
    keepspaces=true,
    %basicstyle={\small\ttfamily},
    basicstyle={\fontsize{8pt}{9pt}\ttfamily},
    numbers=left,
    numberstyle={\fontsize{8pt}{9pt}\ttfamily\color{gray}},
    keywordstyle=\color{webgreen},
    commentstyle=\color{gray},
    stringstyle=\color{Maroon},
    identifierstyle=\color{RoyalBlue},
    breaklines=true,
    breakatwhitespace=true,
    tabsize=3,
    backgroundcolor=\color{codebg},
    captionpos=t,
}

\lstdefinelanguage{ARMAssembler}
    {%
        morekeywords={%
            adc,adcs,%
            add,adds,addw,%
            adr,%
            and,ands,%
            asr,asrs,%
            b,beq,bne,bcs,bcc,bmi,bpl,bvs,bvc,bhi,bls,bge,blt,bgt,ble,bal,bhs,blo,%
            bfc,%
            bfi,%
            bic,bics,%
            bkpt,%
            bl,%
            blx,%
            bx,%
            cbnz,cbz,%
            clrex,%
            clz,%
            cmn,%
            cmp,%
            cps,%
            cpy,%
            csdb,%
            dbg,%
            dmb,%
            dsb,%
            eor,eors,%
            isb,%
            it,itt,ite,ittt,itet,itte,itee,%
            itttt,ittte,ittet,ittee,itett,itete,iteet,iteee,%
            ldc,ldcl,ldc2,ldc2l,%
            ldm,ldmia,ldmfd,%
            ldmdb,ldmea,%
            ldr,%
            ldrb,%
            ldrbt,%
            ldrd,%
            ldrex,%
            ldrexb,%
            ldrexh,%
            ldrh,%
            ldrsb,%
            ldrsbt,%
            ldrsh,%
            ldrsht,%
            ldrt,%
            lsl,lsls,%
            lsr,lsrs,%
            mcr,lcr2,%
            mcrr,mcrr2,%
            mla,%
            mls,%
            mov,movs,movw,movt,%
            mrc,mrc2,%
            mrrc,mrrc2,%
            mrs,%
            msr,%
            mul,muls,%
            mvn,mvns,%
            neg,%
            nop,%
            orn,orns,%
            orr,orrs,%
            pkhbt,pkhtb,%
            pld,%
            pli,%
            pop,%
            pssbb,%
            push,%
            qadd,%
            qadd16,%
            qadd8,%
            qasx,%
            qdadd,%
            qdsub,%
            qsax,%
            qsub,%
            qsub16,%
            qsub8,%
            rbit,%
            rev,%
            rev16,%
            revsh,%
            ror,rors,%
            rrx,%
            rsb,rsbs,%
            sadd16,%
            sadd8,%
            sasx,%
            sbc,sbcs,%
            sbfx,%
            sdiv,%
            sel,sev,%
            shadd16,%
            shadd8,%
            shasx,%
            shsax,%
            shsub16,%
            shsub8,%
            smlabb,smlabt,smlatb,smlatt,%
            smlad,smladx,%
            smlal,%
            smlalbb,smlalbt,smlaltb,smlaltt,%
            smlald,smlaldx,%
            smlawb,smlawt,%
            smlsd,smlsdx,%
            smlsld,smlsldx,%
            smmla,smmlar,%
            smmls,smmlsr,%
            smmul,smmulr,%
            smuad,smuadx,%
            smulbb,smulbt,smultb,smultt,%
            smull,%
            smulwb,smulwt,%
            smusd,smusdx,%
            ssat,%
            ssat16,%
            ssax,%
            ssbb,%
            ssub16,%
            ssub8,%
            stc,stc2,%
            stm,stima,stmea,%
            stmdb,stmfd,%
            str,%
            strb,%
            strbt,%
            strd,%
            strex,%
            strexb,%
            strexh,%
            strh,%
            strht,%
            strt,%
            sub,subs,%
            svc,%
            sxtab,%
            sxtab16,%
            sxtah,%
            sxtb,%
            sxtb16,%
            sxth,%
            tbb,tbh,%
            teq,%
            tst,%
            uadd16,%
            uadd8,%
            uasx,%
            ubfx,%
            udf,%
            udiv,%
            uhadd16,%
            uhadd8,%
            uhasx,%
            uhsax,%
            uhsub16,%
            uhsub8,%
            umaal,%
            umlal,%
            umull,%
            uqadd16,%
            uqadd8,%
            uqasx,%
            uqsax,%
            uqsub16,%
            uqsub8,%
            usad8,%
            usada8,%
            usat,%
            usat16,%
            usax,%
            usub16,%
            usub8,%
            uxtab,%
            uxtab16,%
            uxtah,%
            uxtb,%
            uxtb16,%
            uxth,%
            vabs,%
            vadd,%
            vcmp,vcmpe,%
            vcvta,vcvtn,vcvtp,vcvtm,%
            vcvt,vcvtr,%
            vcvtb,vcvtt,%
            vdiv,%
            vfma,vfms,%
            vldm,%
            vldr,%
            vmaxnm,vminnm,%
            vmla,vmls,%
            vmov,%
            vmrs,%
            vmsr,%
            vneg,%
            vnmla,vnmls,vnmul,%
            vpop,%
            vpush,%
            vrinta,vrintn,vrintp,vrintm,%
            vrintx,%
            vrintz,vrintr,%
            vsel,%
            vsqrt,%
            vstm,%
            vstr,%
            vsub,%
            wfe,%
            wfi,%
            yield},%
        alsoletter={.,0,1,2,3,4,5,6,7,8,9},%
        alsodigit={?},%
        sensitive=false,%
        morestring=[b]",%
        morecomment=[s]{/*}{*/},%
        morecomment=[l]@,%
        morecomment=[l]//,%
    }[keywords,comments,strings]

\lstdefinestyle{zzarmasm}{%
    frame=ltb,
    framerule=0pt,
    aboveskip=3mm,
    belowskip=3mm,
    framextopmargin=.5mm,
    framexbottommargin=.5mm,
    framexleftmargin=1mm,
    framexrightmargin=1.5mm,
    showstringspaces=false,
    language=ARMAssembler,
    columns=fullflexible,
    basicstyle={\fontsize{8pt}{9pt}\ttfamily},
    basewidth=0.5em,
    numbers=left,
    numberstyle=\tiny\color{gray},
    keywordstyle=\color{webgreen},
    commentstyle=\color{gray},
    stringstyle=\color{Maroon},
    identifierstyle=\color{RoyalBlue},
    breaklines=true,
    breakatwhitespace=true,
    tabsize=3,
    backgroundcolor=\color{codebg}
}

\lstdefinestyle{zzc}{%
    frame=ltb,
    framerule=0pt,
    aboveskip=3mm,
    belowskip=3mm,
    framextopmargin=.5mm,
    framexbottommargin=.5mm,
    framexleftmargin=1mm,
    framexrightmargin=1.5mm,
    showstringspaces=false,
    language=C,
    columns=fullflexible,
    basicstyle={\fontsize{8pt}{9pt}\ttfamily},
    basewidth=0.5em,
    numbers=left,
    numberstyle=\tiny\color{gray},
    keywordstyle=\color{webgreen},
    commentstyle=\color{gray},
    stringstyle=\color{Maroon},
    identifierstyle=\color{RoyalBlue},
    breaklines=true,
    breakatwhitespace=true,
    tabsize=3,
    backgroundcolor=\color{codebg}
}

\newenvironment{forbidpagebreak}
    {\par\nobreak\vfil\penalty 0\vfilneg
     \vtop\bgroup}
    {\par\xdef\tpd{\the\prevdepth}\egroup
     \prevdepth=\tpd}

% llncs removes generation of PDF bookmarks, we must add them back
\usepackage{etoolbox}
\makeatletter
\let\llncs@addcontentsline\addcontentsline
\patchcmd{\maketitle}{\addcontentsline}{\llncs@addcontentsline}{}{}
\patchcmd{\maketitle}{\addcontentsline}{\llncs@addcontentsline}{}{}
\patchcmd{\maketitle}{\addcontentsline}{\llncs@addcontentsline}{}{}
\setcounter{tocdepth}{3}
\makeatother
\PassOptionsToPackage{hyphens}{url}
\usepackage[pdftex,bookmarks=true]{hyperref}
\hypersetup{
    bookmarksopen=true,
    bookmarksopenlevel=3,
    bookmarksnumbered=true,
    hidelinks,
    colorlinks,
    linkcolor={red!50!black},
    citecolor={green!50!black},
    urlcolor={blue!80!black}
}
%\hypersetup{hidelinks}
%\usepackage{attachfile2}

\makeatletter
  % Recover some math symbols that were masked by eb-garamond, but do not
  % have replacement definitions.
  \DeclareSymbolFont{ntxletters}{OML}{ntxmi}{m}{it}
  \SetSymbolFont{ntxletters}{bold}{OML}{ntxmi}{b}{it}
  \re@DeclareMathSymbol{\leftharpoonup}{\mathrel}{ntxletters}{"28}
  \re@DeclareMathSymbol{\leftharpoondown}{\mathrel}{ntxletters}{"29}
  \re@DeclareMathSymbol{\rightharpoonup}{\mathrel}{ntxletters}{"2A}
  \re@DeclareMathSymbol{\rightharpoondown}{\mathrel}{ntxletters}{"2B}
  \re@DeclareMathSymbol{\triangleleft}{\mathbin}{ntxletters}{"2F}
  \re@DeclareMathSymbol{\triangleright}{\mathbin}{ntxletters}{"2E}
  \re@DeclareMathSymbol{\partial}{\mathord}{ntxletters}{"40}
  \re@DeclareMathSymbol{\flat}{\mathord}{ntxletters}{"5B}
  \re@DeclareMathSymbol{\natural}{\mathord}{ntxletters}{"5C}
  \re@DeclareMathSymbol{\star}{\mathbin}{ntxletters}{"3F}
  \re@DeclareMathSymbol{\smile}{\mathrel}{ntxletters}{"5E}
  \re@DeclareMathSymbol{\frown}{\mathrel}{ntxletters}{"5F}
  \re@DeclareMathSymbol{\sharp}{\mathord}{ntxletters}{"5D}
  \re@DeclareMathAccent{\vec}{\mathord}{ntxletters}{"7E}

  % Change font for algorithm label.
  \renewcommand\ALG@name{\sffamily\bfseries Algorithm}
\makeatother

% Use the sans-serif fonts (for which bold is properly defined) for
% section headings. Also ensure that subsubsections get a number and
% that the number is displayed (to distinguish them from paragraphs).
\allsectionsfont{\sffamily}
\setcounter{secnumdepth}{3}
\pagestyle{plain}

\makeatletter
\renewenvironment{abstract}{%
      \list{}{\advance\topsep by0.35cm\relax\small
      \leftmargin=1cm
      \labelwidth=\z@
      \listparindent=\z@
      \itemindent\listparindent
      \rightmargin\leftmargin}\item[\hskip\labelsep
                                    \textsf{\textbf{\abstractname}}]}
    {\endlist}
\makeatother
\let\lemma\relax
\spnewtheorem{lemma}{Lemma}{\sffamily\bfseries}{\itshape}

\spnewtheorem{mtheorem}{Theorem}{\sffamily\bfseries}{\itshape}
\spnewtheorem*{mproof}{Proof}{\sffamily\bfseries\itshape}{\rmfamily}

\renewcommand{\algorithmicrequire}{\hbox to3.4em{\textbf{Input:}\hfill}}
\renewcommand{\algorithmicensure}{\hbox to3.4em{\textbf{Output:}\hfill}}

%\newcommand{\GF}{\mathrm{\textit{GF}}}
\newcommand{\GF}{GF}
\newcommand{\QR}{QR}
\newcommand{\bB}{\mathbb{B}}
\newcommand{\bF}{\mathbb{F}}
\newcommand{\bG}{\mathbb{G}}
\newcommand{\bN}{\mathbb{N}}
\newcommand{\bZ}{\mathbb{Z}}
\newcommand{\bR}{\mathbb{R}}
\newcommand{\bQ}{\mathbb{Q}}
\newcommand{\bC}{\mathbb{C}}
\newcommand{\cC}{\mathcal{C}}
\newcommand{\cE}{\mathcal{E}}
\newcommand{\cG}{\mathcal{G}}
\newcommand{\rM}{\text{M}}
\newcommand{\rS}{\text{S}}
\newcommand{\rmb}{\text{m}_\text{b}}
\newcommand{\neutral}{\mathbb{O}}
\newcommand{\vol}{\text{\textsf{vol}}}
\newcommand{\bitlength}{\text{\textsf{len}}}
\newcommand{\MAC}{\text{\textsf{MAC}}}
\newcommand{\MAClen}{\text{\textsf{MAClen}}}
\newcommand{\cc}{\text{\textsf{enc}}}
\newcommand{\cclen}{\text{\textsf{clen}}}
\newcommand{\Setup}{\text{\textsf{Setup}}}
\newcommand{\Eval}{\text{\textsf{Eval}}}
\newcommand{\Verify}{\text{\textsf{Verify}}}
\newcommand{\VDF}{\text{\textsf{VDF}}}
\newcommand{\VDFlen}{\text{\textsf{VDFlen}}}
%\newcommand{\smod}[1]{\,\,\,(\text{mod}^{\pm} #1)}
%\newcommand{\smodnospace}[1]{(\text{mod}^{\pm} #1)}
\newcommand{\sign}{\text{\textsf{sign}}}
\newcommand{\trace}{\text{Tr}}
\newcommand{\qsolve}{\text{QSolve}}
\newcommand{\rev}{\text{rev}}

\newcommand{\ezut}{\textsc{ezut}\xspace}
\newcommand{\xwj}{\textsc{xwj}\xspace}

\newcommand{\eqvm}[3]{#1 \equiv #2 \,[#3]}
\newcommand{\umod}[2]{#1 \,\,\text{mod}^{+}\, #2}
\newcommand{\smod}[2]{#1 \,\,\text{mod}^{\pm}\, #2}
\newcommand{\mont}[1]{\mathcal{M}(#1)}

\newcommand{\norm}[1]{\left\|#1\right\|}
\newcommand{\norminf}[1]{\norm{#1}_{\infty}}
\newcommand{\adj}[1]{{#1}^*}
\newcommand{\cconj}[1]{\overline{#1}}
\newcommand{\inner}[2]{\langle #1, #2 \rangle}

\raggedbottom

\begin{document}

\title{\textsf{A RAM-Efficient Implementation of Falcon}}

\author{Thomas Pornin}
\institute{NCC Group, \email{thomas.pornin@nccgroup.com}}

\maketitle
\noindent\makebox[\textwidth]{7 September, 2026}
%FIXME: adjust date when (re)publishing

\begin{abstract}
We present a RAM-efficient implementation of Falcon: RAM usage has
shrunk to about 11~kB, down from about 31~kB in the previous
implementation of Falcon-512. This code is furthermore faster on Arm
Cortex M4, with average signature generation cost down to 13.45 million
cycles. Optimization techniques include a novel variant of the FFT,
replacement of some floating-point operations with modular integer
computations, delayed addition of input within the Fast Fourier sampling
process, and an alternate signature reassembly process. More than half
of the floating-point operations have been removed. The resulting
implementation is now small enough to allow use in small embedded
systems such as smart cards.
\end{abstract}

% ----------------------------------------------------------------------

\section{Introduction}

Falcon is a post-quantum signature scheme\cite{FalconRound3}; it was
submitted to the NIST post-quantum cryptography standardization
project\cite{NISTPQC} in 2017, and has been selected for becoming the
NIST standard FIPS 206, under the name ``FN-DSA''. Falcon is a
hash-and-sign scheme based on an NTRU lattice. In order to achieve small
signatures and keys while maintaining the expected security level, a
short vector is obtained with an appropriate randomization of Babai's
Nearest Plane algorithm\cite{Bab1986}, as in Klein's
sampler\cite{Kle2000}, except that an NP variant known as Fast Fourier
Orthogonalization is employed, for much improved
performance\cite{DucPre2016}. Falcon's parameter sets use a degree $n =
2^\ell$; the two main choices are $n = 512$ (for NIST's ``level I''
security, similar in strength to AES-128) and $n = 1024$ (for ``level
V'', similar to AES-256).

The verification algorithm of Falcon is very fast, uses little RAM, and
fits in a small code footprint. The key pair generation algorithm is
complicated and large; several previous works have been optimizing that
algorithm\cite{PorPre2019,Por2023-1,Por2025-4}, which now runs with a
large but tolerable CPU cost (e.g. it is still much faster than RSA key
pair generation) and moderate RAM requirements (about 12 or 23~kB,
depending on the degree $n$). The signature generation algorithm,
though, still has significant RAM requirements (about 31~kB for $n =
512$, 61~kB for $n = 1024$), which make it difficult to use on the
smallest embedded systems such as smart cards, where RAM is usually a
very scarce resource. Moreover, the signature generation algorithm, as
specified in~\cite{FalconRound3}, uses floating-point operations, which
are expensive on microcontrollers with no hardware FPU (or whose FPU
does not support the required precision), and whose masking against
side-channel analysis attacks is expected to be doable but
difficult\cite{Ber2026,BerPaiTavBosCol2025,BerTav2025,CheChe2024,CheMinLiuGaoZho2025,KarAys2024}.
This paper aims at reducing both the RAM requirements, and the number of
floating-point operations, in Falcon's signature generation algorithm.

As will be detailed in the rest of the paper, the RAM requirements for
signature generation have been slashed by a factor of 3, down to about
11~kB (for $n = 512$) and 21~kB (for $n = 1024$). More than half of the
floating-point operations have been removed, with the side-effect of
improving performance on an Arm Cortex M4 microcontroller (cost has
dropped to 13.4 million cycles at $n = 512$, down from 17.0 millions).
On large CPUs (x86, aarch64) which have an efficient hardware FPU with
SIMD instructions, this new implementation is moderately slower (cost is
increased by 17\% on a Skylake-class Intel x86, by 27\% on a 64-bit Arm
Cortex A76). This new implementation is available at:
\begin{center}
    \url{https://github.com/pornin/c-fn-dsa-alt}
\end{center}

\paragraph{FN-DSA Standard.}
As these lines are written, the NIST standard FIPS 206 is not published
yet; a draft of that standard has been announced and has been sent into
NIST's approval process at some point around mid-2025, and has since
been stuck there for unclear reasons. There is no available estimate for
when that draft of FIPS 206 will be made public, let alone the final
FIPS 206. For the purposes of this paper, we use as reference the
implementation available at:
\begin{center}
    \url{https://github.com/pornin/c-fn-dsa}
\end{center}
which matches, as of 30 July, 2026, my best guess of what the FIPS 206
draft will contain. This is based on what NIST announced in various
places (workshops, mailing-lists...). Compared to Falcon as submitted,
the main changes, apart from inconsequential endianness choices in
encodings, are the following:
\begin{itemize}
    \item The message to sign is hashed into a message representative $\mu$
    along with some context information and a hash of the public key, as
    is done in ML-DSA\cite{Fips204}. Only $\mu$ is used afterwards.

    \item A new 40-byte random seed is generated for each signing attempt,
    instead of reusing the same seed (this improves provability of the
    security\cite{FouGajGroJanKil2026}).

    \item The infinity norm of a valid signature is enforced to be at
    most 840 (down from 2047).

    \item The public key is in NTT representation.
\end{itemize}
This implementation is what will be called in this paper the
``standard'' implementation. There is no guarantee that when the FIPS
206 draft is published, it will really match this implementation.

\paragraph{Interoperability and Compliance.}
A thorny issue is the question of what is actually a ``compliant''
Falcon/FN-DSA implementation. Right now, with no published standard nor
even draft, there is no real notion of compliance, but NIST's plans seem
to be, for the moment, to require strict reproducibility of test
vectors. An implementation should provide a test API which receives a
source random seed, and everything is deterministically specified from
that seed. The point is to be able to validate that the Gaussian
sampler, which is a rather complicated procedure, has been implemented
properly and does not generate values with an exploitable bias;
verifying strict reproducibility of test vectors goes a long way toward
showing that an implementation at least follows the expected code
structure. Unfortunately, this means that a truly compliant signature
generator MUST perform exactly the same floating-point computations in
the same order as the reference code, because any deviation in the last
few bits of the inputs to the Gaussian sampler will at least
occasionally induce divergences in the rejection sampling process.

The implementation we describe here is \emph{safe} and \emph{fully
interoperable}, but would not be strictly compliant in the sense above:
\begin{itemize}

    \item Safety: all operations are done with enough precision to ensure
    that the security analysis of Falcon still applies (in fact, this
    new implementation has somewhat \emph{better} precision than the
    ``standard'' code, see section~\ref{sec:precision}).

    \item Interoperability: this new implementation uses the same
    formats for public keys, private keys and signatures; signatures
    generated by the ``standard'' implementation are verifiable by our
    new code, and vice versa. \emph{Full} interoperability denotes that
    this also applies to private keys: all private keys generated by one
    implementation are usable by the other (in fact the key pair
    generator code is completely unchanged).

\end{itemize}
Our new implementation will \emph{most of the time} reproduce test vectors,
by using the same private keys, input messages, and source random seeds.
However, divergences will occur with probability about $1/8000$. If NIST
indeed insists on strict reproducibility, then our new code would not be
formally compliant. As explained in section~\ref{sec:hedging}, we
deliberately broke that almost-reproducibility, but it is easy to
reactivate (a one-line code change).

\paragraph{Fixed-Point.}
A recently published implementation uses fixed-point operations instead
of floating-point\cite{AlmFouLacPre2026}. Their implementation is, like
ours, safe and interoperable (it is not always \emph{fully}
interoperable because it enforces some extra requirements on private
keys, so that it cannot use some of the private keys generated by the
``standard'' code). The fixed-point approach should promote performance
on FPU-less hardware, and, more importantly, help with masking
strategies against side-channel attacks. That fixed-point work is
\emph{a priori} complementary with the optimizations described here. A
future work is to merge both implementations, i.e. make the code
presented here use fixed-point operations instead of floating-point.

\paragraph{Main Ideas.}
The main optimization ideas that we use are the following:
\begin{itemize}

    \item We design a new FFT variant, which removes about half of the
    floating-point operations, and furthermore allows halving the
    storage size of some polynomials. This is described in
    section~\ref{sec:fft}. This new FFT might be applicable to other
    schemes than Falcon; it can also be turned into a new NTT that works
    with prime moduli equal to $-1$ modulo $n$, rather than the classic
    $1$ modulo $2n$.

    \item We apply the short representation of polynomials to the
    intermediate node values in the Falcon tree. Moreover, the top-level
    node of the tree is built specially, on-the-fly, with most polynomials
    being computed over integers, using computations modulo $q = 12289$
    and $q' = 18433$; our range analysis shows that the process is
    exact and cannot hit any integer overflow (see section~\ref{sec:tree}).

    \item Only the fractional part of target vectors is used in sampling,
    and it is moreover added later in the process; this improves precision
    and avoids some FFT computations. Similarly, the final output polynomial
    in FFT representation is discarded (and actually not computed at all),
    since the signature is recomputed from the sampled integers; this
    removes the need for any inverse FFT in the whole scheme.
    Section~\ref{sec:genopt} describes these optimizations.

\end{itemize}

% ----------------------------------------------------------------------

\newpage
\section{Falcon Signature Generation}\label{sec:falcon-std}

We recall in this section the Falcon signature generation scheme (and
parts of the key pair generation and verification). Compared with the
Falcon specification\cite{FalconRound3}, we dispatch the operations over
a slightly different structure of sub-algorithms, but the actual
computations are not changed. We also skip the description of some
operations, in particular the Gaussian sampling, that we do not modify.

\subsection{Notations}

\paragraph{Polynomial Ring.}
The main parameter of Falcon is the degree $n = 2^\ell$, for some
integer $\ell \geq 1$. The standard Falcon degree is either $n = 512$
(for ``level I'' security) or $1024$ (for ``level V''). Weak versions,
meant for research purposes and to help development, use lower degrees,
from $n = 4$ to $256$. We consider polynomials modulo the cyclotomic
polynomial $\Phi_n = X^n + 1$. Since $\Phi_n$ is irreducible over $\bQ$,
the ring $\bQ[X]/\Phi_n$ is a field. The same does not hold for
$\bR[X]/\Phi_n$ (for $n\geq 4$) nor for $\bC[X]/\Phi_n$; in the latter,
$\Phi_n$ splits, with roots being $\zeta_{n,j} = e^{i(2j+1)\pi/n}$ for
integers $j = 0$ to $n - 1$.

A polynomial $a \in \bC[X]/\Phi_n$ can be expressed in \emph{coefficient
representation}, i.e. as a sequence of $n$ elements of $\bC$ denoted
$a[j]$ for $j = 0$ to $n - 1$, fulfilling $a = \sum_{j=0}^{n-1} a[j]
X^j$. The \emph{FFT representation} of $a$ is the sequence of $n$ values
$a(\zeta_{n,j})$ for $j = 0$ to $n - 1$. As will be detailed later, the
FFT representation is redundant for the kind of polynomials used in
Falcon, and we omit half or sometimes three quarters of it.

\paragraph{Quotient Ring.}
For $a, b\in\bZ[X]/\Phi_n$ (polynomials with integral coefficients) and
some integer $q > 1$, we can consider equality and operations modulo
$q$: $a = b \bmod q$ if and only if there exists some $w\in
\bZ[X]/\Phi_n$ such that $a = b + qw$. Similarly, we say that $a$ is
invertible modulo $q$ if and only if there exists some $t\in
\bZ[X]/\Phi_n$ such that $at = 1 \bmod q$. Note that this implies that
$a\neq 0$. By extension, we say that $a$ is invertible modulo $1$ if and
only if $a\neq 0$.

Equivalently, we can consider a polynomial ``modulo $q$'' by interpreting
it as being defined over $\bZ_q[X]/\Phi_n$, with $\bZ_q$ being
the ring of integers modulo $q$. If $q$ is prime and is such that
$q = 1 \bmod{2n}$, then multiplications and divisions of polynomials
in $\bZ[X]/\Phi_n$ can be considerably sped up by applying the NTT. In
Falcon, we mainly use $q = 12289$, which matches these conditions. In
the rest of this paper, we will also use the alternate modulus
$q' = 18433$, for which the NTT is also efficient.

\paragraph{Adjoint.}
We define the \emph{(Hermitian) adjoint} of polynomial $a \in
\bC[X]/\Phi_n$ to be the polynomial $\adj{a}$ which is such that
$\adj{a}(\zeta) = \cconj{a(\zeta)}$ for all roots $\zeta$ of $\Phi_n$ in
$\bC$. If the coefficients $a[j]$ are all real, then $\adj{a} = a[0] -
\sum_{j=1}^{n-1} a[n-j] X^j$. For such a real polynomial, we can also
obtain $\adj{a}$ by composition: $\adj{a} = a(1/X)$.

Like complex conjugation, the adjoint is compatible with additions and
multiplications: for polynomials $a$ and $b$, we have $\adj{(a + b)} =
\adj{a} + \adj{b}$, $\adj{(ab)} = \adj{a} \adj{b}$, and $\adj{(a/b)} =
\adj{a}/\adj{b}$.

\paragraph{Inner Product and Norm.}
For polynomials $a, b \in \bR[X]/\Phi_n$, their inner product is:
\begin{equation*}
    \inner{a}{b}
        = \frac{1}{n} \sum_{\Phi_n(\zeta) = 0} a(\zeta) \cconj{b(\zeta)}
        = \sum_{j=0}^{n-1} a[j] b[j]
\end{equation*}
This inner product allows defining the norm of a polynomial $a$ as
$\norm{a} = \sqrt{\inner{a}{a}}$.

We also define the \emph{infinity norm} of a polynomial $a$
as $\norminf{a} = \text{max}_{j=0}^{n-1} |a[j]|$.

The inner product and the norms extend to vectors of polynomials by
computing them coefficient-wise; e.g. $\norm{(a, b)} = \norm{a} + \norm{b}$,
and $\norminf{(a, b)} = \text{max}(\norminf{a}, \norminf{b})$.

\paragraph{Splitting and Merging.}
A polynomial $a \in \bQ[X]/\Phi_n$ can be split into even-indexed and
odd-index coefficients:
\begin{equation*}
    a = a_{\text{even}}(X^2) + X a_{\text{odd}}(X^2)
\end{equation*}
with $a_{\text{even}} = \sum_{j=0}^{n/2-1} a[2j] X^j$ and
$a_{\text{odd}} = \sum_{j=0}^{n/2-1} a[2j+1] X^j$; both
$a_{\text{even}}$ and $a_{\text{odd}}$ are naturally in
$\bQ[X]/\Phi_{n/2}$. The $\text{\textsf{split}}$ and
$\text{\textsf{merge}}$ operations denote splitting and its opposite:
$\text{\textsf{split}}(a) = (a_{\text{even}}, a_{\text{odd}})$, and
$\text{\textsf{merge}}(a_{\text{even}}, a_{\text{odd}}) = a$.

\paragraph{FFT.}
The \emph{Fast Fourier Transform (FFT)} is an efficient way to compute,
for a polynomial $a$, its FFT representation, which is the list of
values $a(\zeta)$ for all roots $\zeta$ of $\Phi_n$ in $\bC$. The gist
of the FFT is related to splitting: if $(a_{\text{even}},
a_{\text{odd}}) = \text{\textsf{split}}(a)$, then, for any root $\zeta$
of $\Phi_n$, $-\zeta$ is also a root of $\Phi_n$, and we have:
\begin{eqnarray*}
    a(\zeta) &=& a_{\text{even}}(\zeta^2) + \zeta a_{\text{odd}}(\zeta^2) \\
    a(-\zeta) &=& a_{\text{even}}(\zeta^2) - \zeta a_{\text{odd}}(\zeta^2)
\end{eqnarray*}
Moreover, $\zeta^2$ is a root of $\Phi_{n/2}$; thus, given a polynomial
$a$, we can split it into $a_{\text{even}}$ and $a_{\text{odd}}$,
recursively apply the FFT on both (in $\Phi_{n/2}$), and then
inexpensively recombine them into the FFT representation of $a$. The FFT
has an asymptotic cost of $O(n \log n)$ operations, which is why it is
computationally attractive. We denote by $\text{\textsf{FFT}}$ and
$\text{\textsf{invFFT}}$ the FFT and its inverse operation,
respectively.

If $a\in\bR[X]/\Phi_n$, then for every root $\zeta$ of $\Phi_n$, its
complex conjugate $\cconj{\zeta}$ is also a root of $\Phi_n$, and
$a(\cconj{\zeta}) = \cconj{a(\zeta)}$ (the real coefficients of $a$ are
unchanged by complex conjugation); thus, we can omit half of the
elements of the FFT representation of $a$, since they can always be
recomputed by complex conjugation.

The FFT representation is compatible with operations on polynomials:
addition, subtraction, multiplication and division of polynomials can be
performed coefficient-wise on the FFT representations of the
polynomials. Unfortunately, the $\zeta$ roots are complex numbers with
irrational real and imaginary parts; this means that in practice, the
FFT representation of a polynomial $a$ must use an approximate
representation, usually some floating-point or fixed-point format.

\paragraph{Vectors and Matrices.}
Some operations use vectors of polynomials, and matrices whose
coefficients are polynomials. They will be denoted with a bold font,
e.g. $\mathbf{t} = (t_0, t_1)$. Vectors use the row convention, i.e.
the vector coordinates are shown horizontally; when applying a linear
operation $\mathbf{B}$ to a vector, the matrix is on the right of
the operand, e.g. $\mathbf{t} \mathbf{B}$.

\subsection{Key Pair Generation}

\paragraph{Base Polynomials and the NTRU Equation.}
Key pair generation consists in sampling two polynomials $f$ and $g$
from $\bZ[X]/\Phi_n$ with a given Gaussian distribution for each
coefficient. The key pair generation algorithm is supposed to enforce
the following rules:
\begin{itemize}

    \item Coefficients of $f$ and $g$ have a limited range, since they
    are encoded in the private key over a limited number of bits: $\norminf{f}
    \leq \delta_n$ and $\norminf{g} \leq \delta_n$ for a value
    $\delta_n$ which depends on $n$: $\delta_{1024} = 16$, $\delta_{512}
    = \delta_{256} = 32$, $\delta_{128} = \delta_{64} = 64$, and
    $\delta_n = 128$ for $n\leq 32$.

    \item $f$ is invertible modulo $q$.

    \item To allow generation of short signatures, some norms are limited:
    \begin{equation*}
        \text{max}\left(\norm{(g, -f)},
        \norm{\left(\frac{q\adj{f}}{f\adj{f} + g\adj{g}},
        \frac{q\adj{g}}{f\adj{f} + g\adj{g}}\right)} \right) \leq
        1.17\sqrt{q}
    \end{equation*}

    \item There exist polynomials $F, G\in\bZ[X]/\Phi_n$ such that $fG -
    gF = q$ (this is called the \emph{NTRU equation}). Moreover,
    $\norminf{F}\leq 127$ and $\norminf{G}\leq 127$. When there is a
    solution $(F, G)$ to the NTRU equation, then there are an infinity
    of solutions (which are all the pairs $(F + kf, G + kg)$ for any
    $k\in \bZ[X]/\Phi_n$), and in general several will match the
    condition on the infinity norms; the key pair generation algorithm
    chooses one solution, whose $F$ is part of the encoding of the
    private key into bytes.

\end{itemize}

$G$ is not part of the encoding of the private key, because it can be
reconstructed from the NTRU equation: given $f$, $g$ and $F$, we have
$G = (q + gF)/f$. Since the division by $f$ is exact over the integers,
$G$ can be evaluated modulo $q$, in which case the equation simplifies
into $G = gF/f \bmod q$; the division can be done coefficient-wise in
NTT representation modulo $q$. Since $\norminf{G} \leq 127 < q/2$, the
result modulo $q$ can be converted to plain integers, yielding the
actual $G$.

There are several known methods to compute $(F, G)$ from $(f, g)$. The
recommended method, which minimizes both RAM and CPU
costs\cite{PorPre2019}, requires that the resultants of $f$ and $g$ with
$\Phi_n$ are odd, coprime integers. Any key pair generation
implementation is free to reject some $(f, g)$ pairs for which it cannot
derive a solution $(F, G)$ even if, mathematically, such a solution
exists. When a candidate $(f, g)$ pair is rejected, a new pair is
sampled (\emph{both} polynomials are resampled; the sampler for $f$ and
$g$ is fully defined by NIST and internally uses SHAKE256).

$f$, $g$, $F$ and $G$ collectively make up the secret NTRU lattice basis:
\begin{equation*}
    \renewcommand*{\arraystretch}{1.3}
    \setlength\arraycolsep{0.3em}
    \mathbf{B} = \begin{bmatrix}
        g & -f \\
        G & -F
    \end{bmatrix}
\end{equation*}
The inverse of this matrix is:
\begin{equation*}
    \renewcommand*{\arraystretch}{1.3}
    \setlength\arraycolsep{0.3em}
    \mathbf{B}^{-1} = \frac{1}{q} \begin{bmatrix}
        -F & f \\
        -G & g
    \end{bmatrix}
\end{equation*}
The same lattice can be generated by the public matrix:
\begin{equation*}
    \renewcommand*{\arraystretch}{1.3}
    \setlength\arraycolsep{0.3em}
    \mathbf{Q} = \begin{bmatrix}
        h & -1 \\
        q & 0 
    \end{bmatrix}
\end{equation*}
with $h = g/f \bmod q$. The Falcon public key is an encoding of $h$.

\paragraph{Falcon Tree.}
The \emph{Falcon tree} is the result of preprocessing a private key
$(f, g, F, G)$ into a number of polynomials which are nominally in
$\bQ[X]/\Phi_n$, but are usually expressed in FFT representation. These
polynomials are organized in a binary tree structure:
\begin{itemize}

    \item Inner nodes have a \textsf{value}, which is a polynomial, and
    two children (labelled \textsf{left} and \textsf{right}) which are
    themselves trees.

    \item Leaf nodes are real numbers.

    \item A tree is represented by its root node, which is either an
    inner node, or a leaf value.

    \item The full tree starts with a root whose polynomial value is in
    $\bQ[X]/\Phi_n$; the two children of a tree node that contains a
    polynomial in $\bQ[X]/\Phi_m$ themselves have values that are
    polynomials in $\bQ[X]/\Phi_{m/2}$. The degree is thus halved at
    each level; when an inner node holds a \textsf{value} expressed in
    $\bQ[X]/\Phi_2$, its two children are leaves, i.e. real numbers.

\end{itemize}

The algorithm \textsf{FalconTreeTop} (algorithm~\ref{alg:FalconTreeTop})
computes the Falcon tree; it internally uses \textsf{FalconTreeInner}
(algorithm~\ref{alg:FalconTreeInner}), which is recursive.

\begin{algorithm}
    \caption{\ \ \textsf{FalconTreeTop}\label{alg:FalconTreeTop}}
    \begin{algorithmic}[1]
        \Require{$f, g, F, G \in \bQ[X]/\Phi_n$}
        \Ensure{Binary tree \textsf{T}}
        \State{$a \leftarrow f\adj{f} + g\adj{g}$
            \Comment{$a\in\bQ[X]/\Phi_n$ and $\adj{a} = a$}}
        \State{$b \leftarrow F\adj{f} + G\adj{g}$}
        \State{$c \leftarrow F\adj{F} + G\adj{G}$}
        \State{$d \leftarrow c - b\adj{b}/a$
                \Comment{$d\in\bQ[X]/\Phi_n$ and $\adj{d} = d$}}
        \State{$L \leftarrow b/a$}
        \State{$\text{\textsf{T.value}} \leftarrow L$}
        \State{$\text{\textsf{T.left}}
            \leftarrow \text{\textsf{FalconTreeInner}}(a)$}
        \State{$\text{\textsf{T.right}}
            \leftarrow \text{\textsf{FalconTreeInner}}(d)$}
        \State{\Return{\textsf{T}}}
    \end{algorithmic}
\end{algorithm}

\begin{algorithm}
    \caption{\ \ \textsf{FalconTreeInner}\label{alg:FalconTreeInner}}
    \begin{algorithmic}[1]
        \Require{$a \in \bQ[X]/\Phi_m$ such that $a = \adj{a}$}
        \Ensure{Binary tree \textsf{T}}
        \State{$(b, c) \leftarrow \text{\textsf{split}}(a)$
            \Comment{$b, c\in\bQ[X]/\Phi_{m/2}$; also: $\adj{b} = b$}}
        \State{$L \leftarrow \adj{c}/b$}
        \State{$d \leftarrow b - c\adj{c}/b$ \Comment{$\adj{d} = d$}}
        \State{$\text{\textsf{T.value}} \leftarrow L$}
        \If{$m \geq 8$}
            \State{$\text{\textsf{T.left}}
                \leftarrow \text{\textsf{FalconTreeInner}}(b)$}
            \State{$\text{\textsf{T.right}}
                \leftarrow \text{\textsf{FalconTreeInner}}(d)$}
        \Else
            \State{$\text{\textsf{T.left}} \leftarrow b(0)$}
            \State{$\text{\textsf{T.right}} \leftarrow d(0)$}
        \EndIf
        \State{\Return{\textsf{T}}}
    \end{algorithmic}
\end{algorithm}

The Falcon specification describes the tree building with a single
function (called $\text{\textsf{ffLDL}}^\star$) which takes over the
roles of both \textsf{FalconTreeTop} and \textsf{FalconTreeInner}. Here,
we use two separate functions to highlight some properties that we will
use later on, in particular that the input of \textsf{FalconTreeInner}
is a single polynomial which is futhermore equal to its adjoint.

In the Falcon specification, once the tree has been computed, its leaves
are normalized: if a leaf value is $x$, then it is replaced with
$\sigma_n/\sqrt{x}$ for a constant $\sigma_n$ defined as:
\begin{equation*}
    \sigma_n = 1.17 \sqrt{\frac{q \log\left(4n\left(1 + 2^{32}\sqrt{\text{max}(2, n/4)}\right)\right)}{2\pi^2}}
\end{equation*}
In the description shown here, the normalization step is performed right
before using the leaf in a call to \textsf{SamplerZ}, even though it could
be done as part of the tree building.

Computing the Falcon tree involves splitting polynomials, and computing
some arithmetic operations on these polynomials. In plain coefficient
representation, splitting a polynomial is trivial, but polynomial
divisions are very much not so. It is possible to compute a division in
$\bQ[X]/\Phi_n$ by performing an extended Euclidean GCD between the
divisor and $\Phi_n$. Another method, which is applicable here thanks to
the structure of $\Phi_n$, is to perform a descent in a tower of fields:
for a polynomial $b\in\bQ[X]/\Phi_n$, we can define the Galois conjugate
of $b$ as $b^\times(X) = b(-X)$ (i.e. we simply negate the odd-indexed
coefficients), and then compute $a/b = (ab^\times)/(bb^\times)$. The
product $bb^\times$ happens to have only even-indexed coefficients (the
odd-indexed coefficients are null), hence it can be expressed as
$bb^\times = b'(X^2)$ for a polynomial $b' \in\bQ[X]/\Phi_{n/2}$. Using
this method recursively on $b'$, we finally transform the original
fraction $a/b$ into $c/r$ with $r$ being a constant polynomial; at that
point, the division can be applied on the coefficients of $c$ one by
one. As detailed in \cite{PorPre2019}, the value $r$ is really the
resultant of $b$ with $\Phi_n$.

Regardless of the division method used, the resulting polynomial
coefficients will have very large denominators, in the thousands of
bits, implying large RAM usage and poor performance. To make the tree
building tolerably efficient, the Falcon specification uses the FFT: the
source polynomials $f$, $g$, $F$ and $G$ are converted to the FFT
representation (as $n/2$ complex values each), and all operations are
performed in the FFT domain. Addition, subtraction, multiplication and
division of polynomials are then done coefficient-wise. Splitting in
the FFT domain is denoted \textsf{splitfft} in the Falcon specification,
and can be done efficiently: for any root $\zeta$ of $\Phi_n$, $\zeta^2$
is a root of $\Phi_{n/2}$, and:
\begin{eqnarray*}
    a_{\text{even}}(\zeta^2) &=& (a(\zeta) + a(-\zeta))/2 \\
    a_{\text{odd}}(\zeta^2) &=& (a(\zeta) - a(-\zeta))/(2\zeta)
\end{eqnarray*}

The Falcon tree depends only on the private key (the $f$, $g$, $F$ and
$G$ polynomials) and, as such, can be precomputed, though this is not
mandatory; it is also possible to recompute it at signing time.
Signature generation involves a depth-first walk of the tree in
right-to-left order; this can be combined with a dynamic generation that
only needs to temporarily store the polynomials in the path from the
root to the current leaf, thus not needing the RAM space that would be
required for storing the whole tree.

\subsection{Signature Generation}

Signature of a message $M$ first requires pre-processing $M$ along with
some contextual information (e.g. if $M$ is really a hashed value, then
an identifier for the used hash function is included) and a hash of the
public key; the output is a 64-byte string $\mu$. Then, a random 40-byte
seed $r$ is generated, and hashed along with $\mu$ to feed a rejection
sampling process that outputs a polynomial $c\in\bZ_q[X]/\Phi_n$. The
40-byte seed will be included in the signature, so that verifiers may
recompute $c$. If the signature generation fails (it can happen, with a
relatively low probability), a new 40-byte seed is
generated\footnote{This is one of the differences between the original
Falcon\cite{FalconRound3} and the presumed future FIPS 206; in the
original Falcon, $r$ was reused in that case.}, but $\mu$ is reused. We
omit here the details of the generation of $c$.

We suppose for now that the private key was decoded into $f$, $g$ and
$F$, then $G$ was recomputed using the NTRU equation modulo $q$, and
the Falcon tree $\text{\textsf{T}}$ was built. Given $c$, the signature
generation proceeds as follows:
\begin{enumerate}
    \item Compute $\mathbf{t} = (c, 0) \mathbf{B}^{-1} = (-cF/q, cf/q)$.

    \item Sample a vector $\mathbf{z}$ of integral polynomials
    through the \textsf{ffSampling} algorithm.

    \item Compute $\mathbf{s} = (\mathbf{t} - \mathbf{z}) \mathbf{B}$.

    \item If $\norm{\mathbf{s}}$ and $\norminf{\mathbf{s}}$ are low
    enough, then try to encode the second element of $\mathbf{s}$
    into bytes; if the encoding succeeds, then this is the signature.
    Otherwise, the process starts again with a new random 40-byte seed
    $r$ and thus a new polynomial $c$.

\end{enumerate}

The intuition behind this process is the following: the signature is a
vector $\mathbf{e}$ which is part of the lattice defined by the public
basis $\mathbf{Q}$ (for which $\mathbf{B}$ is also a basis) and which is
\emph{close} to the target vector $(c, 0)$. The lattice consists of all
vectors $\mathbf{k} \mathbf{B}$ for any integral polynomial
$\mathbf{k}$; the vector $\mathbf{t}$ is such that $\mathbf{t}
\mathbf{B}$ is \emph{equal} to $(c, 0)$, but $\mathbf{t}$ is not
integral ($\mathbf{B}^{-1}$ starts with a $1/q$ factor). The
\textsf{ffSampling} procedure finds a vector $\mathbf{z}$ which is
integral such that $\mathbf{z} \mathbf{B}$ is close to $\mathbf{t}
\mathbf{B}$, at which point we can set $\mathbf{e} = \mathbf{z}
\mathbf{B}$. The signature itself is then an encoding of $\mathbf{s} =
(c, 0) - \mathbf{e} = (\mathbf{t} - \mathbf{z}) \mathbf{B}$, which has
low norm. For appropriate security, $\mathbf{e}$ must not be \emph{too
close} to the target $(c, 0)$, since that would otherwise leak too much
information about the secret basis $\mathbf{B}$; the \textsf{ffSampling}
procedure is a variant of Klein's sampler\cite{Kle2000}, which is itself
a randomized variant of Babai's Nearest Plane algorithm\cite{Bab1986}.
Klein's sampler implementation has a large cost (in $O(n^2)$);
\textsf{ffSampling} organizes it with a recursive splitting process
which benefits from the FFT representation, for a $O(n\log n)$
cost\cite{DucPre2016}.

\textsf{ffSampling} (algorithm~\ref{alg:ffSampling}) uses the Falcon
tree; namely, it receives as parameters a degree $m$, a target vector
$\mathbf{t}$ of polynomials in $\bQ[X]/\Phi_m$, and a Falcon tree whose
root node holds a value which is itself an element of $\bQ[X]/\Phi_m$.
The process is recursive, the degree $m$ being halved whenever going one
level deeper in the recursion; when $m = 1$, the tree is reduced to a
single real value, which is used as one of the parameters of
\textsf{SamplerZ}.

\begin{algorithm}
    \caption{\ \ \textsf{ffSampling}\label{alg:ffSampling}}
    \begin{algorithmic}[1]
        \Require{degree $m$,
            target $\mathbf{t} = (t_0, t_1) \in (\bQ[X]/\Phi_m)^2$,
            a Falcon tree \textsf{T} whose root has degree $m$}
        \Ensure{$\mathbf{z} = (z_0, z_1) \in (\bZ[X]/\Phi_m)^2$}
        \If{$m = 1$}
            \State{$\sigma' \leftarrow \sigma_n/\text{\textsf{T}}$\Comment{At
                degree 1, the tree is a single real number}}.
            \State{$z_0 \leftarrow \text{\textsf{SamplerZ}}(t_0, \sigma')$}
            \State{$z_1 \leftarrow \text{\textsf{SamplerZ}}(t_1, \sigma')$}
            \State{\Return{$(z_0, z_1)$}}
        \EndIf
        \State{$\mathbf{v} \leftarrow \text{\textsf{split}}(t_1)$}
        \State{$\mathbf{y} \leftarrow \text{\textsf{ffSampling}}(m/2,
            \mathbf{v}, \text{\textsf{T.right}})$}
        \State{$z_1 \leftarrow \text{\textsf{merge}}(\mathbf{y})$}
        \State{$L \leftarrow \text{\textsf{T.value}}$}
        \State{$t'_0 \leftarrow t_0 + (t_1 - z_1)L$}
        \State{$\mathbf{u} \leftarrow \text{\textsf{split}}(t'_0)$}
        \State{$\mathbf{w} \leftarrow \text{\textsf{ffSampling}}(m/2,
            \mathbf{u}, \text{\textsf{T.left}})$}
        \State{$z_0 \leftarrow \text{\textsf{merge}}(\mathbf{w})$}
        \State{\Return{$(z_0, z_1)$}}
    \end{algorithmic}
\end{algorithm}

The \textsf{SamplerZ} function takes as inputs a target real number $x$
and a width $\sigma'$, and it returns an integer $z$ sampled from a
Gaussian distribution centred on $x$ and with standard deviation
$\sigma'$. \textsf{SamplerZ} internally consumes random bits from a
random source. We will not detail here how \textsf{SamplerZ} works,
except to express an important characteristic, which is that if
$\text{\textsf{SamplerZ}}(x, \sigma')$ returns the integer $z$ for a
given state of its internal random source, then, for any integer $j$,
$\text{\textsf{SamplerZ}}(x + j, \sigma')$ should return $z + j$ for the
same internal random source state, and consume exactly as many bits from
that source. In other words, \textsf{SamplerZ} first splits the input
$x$ into an integer $x_i$ and a fractional part $x_f$ such that $x = x_i
+ x_f$ and $0\leq x_f < 1$; it then performs the sampling using $x_f$
alone, and finally adds back $x_i$ to the sampled integer.

\subsection{Signature Verification}

Signature verification works over the message $M$, the public lattice
basis $\mathbf{Q}$, and the signature, which nominally includes $r$ and
the vector $\mathbf{s} = (s_0, s_1)$. The verifier computes the message
representative $\mu$, then (using $r$) the target polynomial $c$. The
signature is deemed valid if all of the following hold:
\begin{eqnarray*}
    c &=& s_0 + s_1 h \mod q \\
    \norm{\mathbf{s}} &\leq& \beta \\
    \norminf{\mathbf{s}} &\leq& \gamma_\infty
\end{eqnarray*}
The thresholds for the last two conditions are scheme parameters:
$\beta = 1.1 \sigma_n \sqrt{2n}$, and $\gamma_\infty = 840$. The first
condition really expresses that $(c, 0) - \mathbf{s}$ is part of the
lattice generated by the public matrix $\mathbf{Q}$. Since this equation
allows recomputing $s_0$ from $s_1$ (using $s_0 = c - s_1 h \bmod q$),
only $s_1$ and the random seed $r$ need to be actually encoded in the
signature. As for the public matrix $\mathbf{Q}$, its contents can be
conveyed by encoding only the public polynomial $h$, whose coefficients
are all integers in the $[0, q-1]$ range.

% ----------------------------------------------------------------------

\newpage
\section{Optimizing the FFT}\label{sec:fft}

In the Falcon specification, the FFT is used throughout the process in
order to speed up operations and in particular polynomial divisions.
Polynomials $f$, $g$, $F$, $G$ and $c$ are all converted to FFT
representation; the computations of $\mathbf{t} = (-cF/q, cf/q)$ and all
\textsf{split}, \textsf{merge} and arithmetic operations on polynomials
in \textsf{ffSampling} are then performed on FFT representations. At the
deepest recursion level, polynomials $t_0$ and $t_1$ are modulo $\Phi_1
= X + 1$, i.e. they have degree zero, and are equal to their FFT
representation, so that their values can be used directly in
\textsf{SamplerZ}. Finally, the output of \textsf{ffSampling} can be
used to compute $\mathbf{s}$, still working with FFT representations,
and only at the end is $\mathbf{s}$ reconverted back to coefficient
representation through the inverse FFT. In total, following the Falcon
specification, five FFTs and two inverse FFTs are performed on
polynomials. Making these operations more efficient thus seems important
for performance.

Here, we primarily target small embedded systems, such as
microcontrollers using an Arm Cortex M4 core. Such systems do not offer
hardware support for floating-point operations with the precision needed
by Falcon; these operations must then be emulated with integer code,
which should furthermore avoid leaking secret information through
side-channels, and thus should at least operate in a constant-time way.
The cost of floating-point operations thus dominates that of the FFT,
and the primary goal is to reduce the number of such operations. Our
assembly-optimized implementations have the following costs:
\begin{itemize}

    \item Addition: 72 cycles

    \item Combined addition-subtraction: 87 cycles

    \item Multiplication: 36 cycles

\end{itemize}
A combined addition-subtraction is obtaining both $x + y$ and $x - y$
for given inputs $x$ and $y$. Addition is counter-intuitively more
expensive than multiplication because it involves shifting operands to
make them share the same exponent, and that shift is expensive,
especially since the shift amount may or may not exceed the size of a
register, and constant-time discipline forces us to always cover all
cases. Negation, and scaling by a power of two (in particular, dividing
a value by 2), are inexpensive, so that the cost of a subtraction can be
considered to be the same as that of an addition (in all of the
subsequent text, we will loosely talk of ``additions'' to designate both
additions and subtractions, since they have the same cost). These
figures are in the context of our Falcon implementation, which uses the
IEEE 754 ``binary64'' format, and for which it can be proven that NaNs,
infinites and denormals never occur, neither as source operands nor as
outputs; the only special values are zeros.

We will here estimate costs in numbers of additions (\textsc{a}) (which
can be subtractions in practice), combined addition-subtractions
(\textsc{as}) and multiplications (\textsc{m}).

\subsection{Classic FFT}

We call ``classic FFT'' the algorithm that applies the ``Cooley-Tukey
butterfly'' repeatedly\cite{CooTuk1965}\footnote{The algorithm was
actually invented by Gauss more than a century before\cite{Gau1866}.},
which really is about computing $a(\zeta)$ and $a(-\zeta)$, for roots
$\zeta$ of $\Phi_m$, from the values $a_{\text{even}}(\zeta^2)$ and
$a_{\text{odd}}(\zeta^2)$, which are the FFT representations of
$a_{\text{even}}$ and $a_{\text{odd}}$ in the ring $\bR[X]/\Phi_{n/2}$.
The corresponding inverse FFT simply reverts all operations using the
reverse of the Cooley-Tukey butterfly, usually known as the
``Gentleman-Sande butterfly''\cite{GenSan1966}. Practical
implementations work in-place, to avoid extra memory allocation and
unneeded data movement, which has the side-effect of returning the FFT
representation in so-called ``bit-reversal order''. We define the
\textsf{rev} function:
\begin{equation*}
    \text{\textsf{rev}}\left(\alpha, \sum_{j=0}^{\alpha-1} x_j 2^j\right) =
        \sum_{j=0}^{\alpha-1} x_j 2^{\alpha - 1 - j}
\end{equation*}
with values $x_j \in [0, 1]$ (i.e. the $x_j$ values are the digits of
the representation of an integer from $[0, 2^\alpha - 1]$ in radix two, and
the \textsf{rev} function reverses the order of that bit sequence). The
classic FFT and its inverse are shown in algorithms \textsf{FFT-Classic}
(algorithm~\ref{alg:fft-classic}) and \textsf{invFFT-classic}
(algorithm~\ref{alg:ifft-classic}).

\begin{algorithm}
    \caption{\ \ \textsf{FFT-classic}\label{alg:fft-classic}}
    \begin{algorithmic}[1]
        \Require{$a = \sum_{j=0}^{n-1} a[j] X^j \in \bR[X]/\Phi_n$,
            for $n = 2^\ell$}
        \Ensure{FFT representation of $a$ in bit-reversal order}
        \For{$\alpha = 1$ \textbf{to} $\ell - 1$}
            \State{$t \leftarrow n/2^{\alpha+1}$}
            \State{$m \leftarrow 2^\alpha$}
            \For{$j = 0$ \textbf{to} $m/2 - 1$}
                \State{$s_r \leftarrow Z_r[m + j]$\Comment{
                    $Z_r[\text{\textsf{rev}}(\ell, j)] = \cos(\pi j/n)$}}
                \State{$s_i \leftarrow Z_i[m + j]$\Comment{
                    $Z_i[\text{\textsf{rev}}(\ell, j)] = \sin(\pi j/n)$}}
                \For{$k = 0$ \textbf{to} $t/2 - 1$}
                    \State{$k_0 \leftarrow 2^{\ell - \alpha} j + k$}
                    \State{$k_1 \leftarrow k_0 + t$}
                    \State{$(x_r, x_i) \leftarrow (a[k_0], a[k_0 + n/2])$}
                    \State{$(y_r, y_i) \leftarrow (a[k_1], a[k_1 + n/2])$}
                    \State{$z_r \leftarrow y_r s_r - y_i s_i$}
                    \State{$z_i \leftarrow y_r s_i + y_i s_r$}
                    \State{$(a[j_0], a[j_0 + n/2])
                        \leftarrow (x_r + z_r, x_i + z_i)$}
                    \State{$(a[j_1], a[j_1 + n/2])
                        \leftarrow (x_r - z_r, x_i - z_i)$}
                \EndFor
            \EndFor
        \EndFor
    \end{algorithmic}
\end{algorithm}

\begin{algorithm}
    \caption{\ \ \textsf{invFFT-classic}\label{alg:ifft-classic}}
    \begin{algorithmic}[1]
        \Require{FFT representation of $a$ in bit-reversal order}
        \Require{$a = \sum_{j=0}^{n-1} a[j] X^j \in \bR[X]/\Phi_n$}
        \For{$\alpha = \ell - 1$ \textbf{down to} $1$}
            \State{$t \leftarrow 2^{\alpha - 1}$}
            \State{$m \leftarrow n/2^\alpha$}
            \For{$j = 0$ \textbf{to} $m/2 - 1$}
                \State{$s_r \leftarrow Z_r[m + j]$
                    \Comment{Same $Z_r[]$ array as in \textsf{FFT-classic}.}}
                \State{$s_i \leftarrow -Z_i[m + j]$
                    \Comment{Same $Z_i[]$ array as in \textsf{FFT-classic}.}}
                \For{$k = 0$ \textbf{to} $t - 1$}
                    \State{$k_0 \leftarrow 2^\alpha j + k$}
                    \State{$k_1 \leftarrow k_0 + t$}
                    \State{$(x_r, x_i) \leftarrow (a[k_0], a[k_0 + n/2])$}
                    \State{$(y_r, y_i) \leftarrow (a[k_1], a[k_1 + n/2])$}
                    \State{$(a[j_0], a[j_0 + n/2])
                        \leftarrow ((x_r + y_r)/2, (x_i + y_i)/2)$}
                    \State{$(z_r, z_i)
                        \leftarrow ((x_r - y_r)/2, (x_i - y_i)/2)$}
                    \State{$a[k_1] \leftarrow z_r s_r - z_i s_i$}
                    \State{$a[k_1 + n/2] \leftarrow z_r s_i + z_i s_r$}
                \EndFor
            \EndFor
        \EndFor
    \end{algorithmic}
\end{algorithm}

The following points are noteworthy:
\begin{itemize}

    \item Since the source polynomials are from $\bR[X]/\Phi_n$, the FFT
    representation only returns $n/2$ complex values (the other $n/2$
    values are implicit since they can be obtained through conjugation).
    \textsf{FFT-classic} returns $a(\zeta_{n,2j})$ for $j = 0$ to $n/2-1$,
    with $\zeta_{n,2j} = e^{i(4j+1)\pi/n}$.

    \item The FFT coefficients are in bit-reversal order, with the real
    and imaginary parts separate; namely:
    \begin{equation*}
        a(\zeta_{n,\text{\textsf{rev}}(\ell, j)})
            = {\hat a}[j] + i {\hat a}[j + n/2]
    \end{equation*}
    for $j = 0$ to $n/2 - 1$, and ${\hat a}$ denoting the contents of
    the $a$ array when \textsf{FFT-classic} terminates. This separation
    of the real and imaginery parts implies that the first formal
    iteration of the FFT algorithm does nothing (which is why the
    counter $\alpha$ starts at 1 in algorithm~\ref{alg:fft-classic},
    instead of 0); it also promotes use of SIMD instruction sets since
    memory slots at successive addresses are processed identically.

    \item Constants $Z_r[]$ and $Z_i[]$ take up a total of 16384 bytes
    of code space to support degrees up to $n = 1024$ and using 64-bit
    floating-point values; however, this size can easily be reduced to
    4096 bytes by using the facts that $\sin(x) = \sin(\pi - x)$ and
    $\cos(x) = \sin(\pi/2 - x)$, so that only the values of $\sin(\pi
    j/n)$ for $j = 0$ to $n/2 - 1$ really need to be stored.

    \item The description of \textsf{FFT-classic} details the operations
    on real and imaginary parts separately, and not using an abstract
    notation with complex numbers, to highlight the cost of the
    multiplication by $\zeta$ (which is $s_r + i s_i$ in the algorithm):
    such an operation requires four floating-point multiplications and
    two floating-point additions\footnote{This could also be implemented
    with three multiplications and five additions, with
    $(y_r + y_i)(s_r + s_i) = y_r s_r + y_i s_i + (y_r s_i + y_i s_r)$,
    and computing $y_r s_r$ and $y_i s_i$ separately; however, this would
    be more expensive since floating-point additions are more expensive
    than floating-point multiplications.}. Our main optimization technique
    aims at avoiding multiplications of complex values.

\end{itemize}

Algorithm \textsf{FFT-classic} involves $\ell - 1$ iterations of the
outer loop; each iteration involves combining the $n$ elements by pairs,
each combination involving two additions, two addition-subtractions, and
four multiplications. In the first iteration, some marginal savings are
possible: $m = 2$, so that $s_r = Z_r[2] = \cos(\pi/4) = \sin(\pi/4) =
Z_i[2] = s_i$; thus, in that first outer iteration, $y_r s_r = y_r s_i$,
$y_i s_i = y_i s_r$, and each pair may use three addition-subtractions
and two multiplications. The total cost of \textsf{FFT-classic} is then
$(n/2)((2\ell - 4)\text{\textsc{a}} + (2\ell - 1)\text{\textsc{as}} +
(4\ell - 6)\text{\textsc{m}})$. Analysis for \textsf{invFFT-classic} is
similar and yields the same cost estimate.

\subsection{FFT of a Self-Adjoint Polynomial}

Let $a \in \bR[X]/\Phi_n$. We define that:
\begin{itemize}

    \item $a$ is \emph{self-adjoint} if and only if: $\adj{a} = a$

    \item $a$ is \emph{$X$-adjoint} if and only if: $\adj{a} = X a$

    \item $a$ is \emph{$1/X$-adjoint} if and only if: $\adj{a} = (1/X) a$

    \item $a$ is \emph{neg-adjoint} if and only if: $\adj{a} = -a$

\end{itemize}

The following remark shows that for all four categories, the full
representation is redundant, and half of the coefficients can be inferred
from the other half:
\begin{remark}\label{remark:adjcoeffs}
For $a \in \bR[X]/\Phi_n$:
\begin{itemize}

    \item If $a$ is self-adjoint then $a[n/2] = 0$, and
    $a[n/2 + j] = -a[n/2 - j]$ for $j = 1$ to $n/2 - 1$.

    \item If $a$ is $X$-adjoint then $a[j] = -a[n - 1 - j]$ for
    $j = 0$ to $n/2 - 1$.

    \item If $a$ is $1/X$-adjoint then $a[1] = a[0]$, and
    $a[n + 1 - j] = -a[j]$ for $j = 2$ to $n/2$.

    \item If $a$ is neg-adjoint then $a[0] = 0$, and
    $a[n/2 + j] = a[n/2 - j]$ for $j = 1$ to $n/2 - 1$.

\end{itemize}
\end{remark}

As previously defined, for any polynomial $a\in \bR[X]/\Phi_n$,
$\adj{a}(\zeta) = \cconj{a(\zeta)}$ for all roots $\zeta$ of $\Phi_n$.
Apart from giving an easy way to transform a polynomial into its adjoint
in FFT representation (just negate the imaginary part of each FFT
coefficient), this implies that if $a$ is self-adjoint, then all its FFT
coefficients are real. In particular, in the representation used by
\textsf{FFT-classic}, elements of index $n/2$ to $n-1$ are the imaginary
parts of the FFT coefficients, and thus they are all equal to zero in a
self-adjoint polynomial, and can be omitted. Moreover, for both
\textsf{FFT-classic} and \textsf{invFFT-classic}, values from the first
half of the array ($a[0]$ to $a[n/2-1]$) are never combined with values
from the second half ($a[n/2]$ to $a[n-1]$). This allows making variants
\textsf{FFT-selfadj-classic} and \textsf{invFFT-selfadj-classic} for
self-adjoint polynomials: they work with half-size arrays ($n/2$
elements instead of $n$) and have half the cost of the full
\textsf{FFT-classic} and \textsf{invFFT-classic}, i.e. $(n/4)((2\ell -
4)\text{\textsc{a}} + (2\ell - 1)\text{\textsc{as}} + (4\ell -
6)\text{\textsc{m}})$.

We can do better. Some easily proven properties are the following:
\begin{remark}\label{remark:split-selfadj}
Let $a \in \bR[X]/\Phi_n$ a self-adjoint polynomial. If
$(a_{\text{even}}, a_{\text{odd}}) = \text{\textsf{split}}(a)$, then
$a_{\text{even}}$ is self-adjoint, and $a_{\text{odd}}$ is $X$-adjoint.
\end{remark}

\begin{remark}
In $\bR[X]/\Phi_n$, $1 + X$ is $1/X$-adjoint, and $X^{n/2}$ is neg-adjoint.
\end{remark}

Let's consider $a \in \bR[X]/\Phi_n$ a self-adjoint polynomial, and
$(a_{\text{even}}, a_{\text{odd}}) = \textsf{split}(a)$. Define
$b_{\text{odd}} = (1 + X) a_{\text{odd}}$ (this is multiplication in
$\bR[X]/\Phi_{n/2}$, since the output of \text{\textsf{split}} consists
of half-degree polynomials). Then:
\begin{equation*}
    \adj{b}_{\text{odd}} = \adj{(1 + X)} \adj{a}_{\text{odd}}
        = (1/X) X (1 + X) a_{\text{odd}} = b_{\text{odd}}
\end{equation*}
Thus, $b_{\text{odd}}$ is self-adjoint. When applying polynomial $a$ to
the roots $\zeta$ of $\Phi_n$, we obtain:
\begin{eqnarray*}
    a(\zeta) &=& a_{\text{even}}(\zeta^2) + \zeta a_{\text{odd}}(\zeta^2) \\
             &=& a_{\text{even}}(\zeta^2)
                 + \frac{\zeta}{1 + \zeta^2} b_{\text{odd}}(\zeta^2)
\end{eqnarray*}
Note that $a$, $a_{\text{even}}$ and $b_{\text{odd}}$ are all
self-adjoint, and thus $a(\zeta)$, $a_{\text{even}}(\zeta^2)$ and
$b_{\text{odd}}(\zeta^2)$ are real numbers. It follows that $\zeta/(1 +
\zeta^2) \in \bR$, and, crucially, the multiplication of that value with
$b_{\text{odd}}(\zeta^2)$ is between two reals, not two complex values,
thus involving a single floating-point multiplication and no addition at
all. Indeed, we have:
\begin{equation*}
    \frac{\zeta_{n,j}}{1 + \zeta_{n,j}^2}
        = \frac{e^{i(2j+1)\pi/n}}{1 + e^{i(2(2j+1))\pi/n}}
        = \frac{1}{e^{-i(2j+1)\pi/n} + e^{i(2j+1)\pi/n}}
        = \frac{1}{2\cos ((2j+1)\pi/n)}
\end{equation*}
This yields an FFT algorithm for self-adjoint polynomials which works
in two steps:
\begin{enumerate}
    \item Split the source polynomial into even-indexed and odd-indexed
    coefficients, then multiply by $1 + X$ the odd-indexed half, and
    apply the same process recursively on both halves.
    \item Recombine the polynomials by pairs, starting with degree 1
    and upward, using the equations above and the combination values
    $\zeta/(1 + \zeta^2)$.
\end{enumerate}
Algorithm \textsf{FFT-selfadj-new} (algorithm~\ref{alg:fft-selfadj-new})
implements this algorithm. The constants $R[j]$ are such that:
\begin{equation*}
    R[2^{\alpha - 2} + \text{\textsf{rev}}(\alpha - 2, j)]
        = \frac{1}{2 \cos ((4 j + 1)\pi/2^\alpha)}
\end{equation*}
for all $j = 0$ to $2^{\alpha - 2} - 1$. $R[0]$ is not used, and $R[1] =
\sqrt{2}/2$. To support polynomials of degree $n = 2^\ell$, the array
$R[]$ must be computed for up to $\alpha = \ell$, which means a total of
$2^{\ell - 1} = n/2$ elements. For Falcon, we need to support up to $n =
1024$ and we use 64-bit floating-point values, so that the footprint of
the constant $R[]$ array is 4096 bytes.

\begin{algorithm}
    \caption{\ \ \textsf{FFT-selfadj-new}\label{alg:fft-selfadj-new}}
    \begin{algorithmic}[1]
        \Require{$a[0]$ to $a[n/2 - 1]$ for $a\in\bR[X]/\Phi_n$ such
            that $n = 2^\ell \geq 4$ and $\adj{a} = a$}
        \Ensure{FFT representation of $a$ in bit-reversal order (half-size)}
        \For{$\alpha = 0$ \textbf{to} $\ell - 3$}
                \Comment{Loop skipped if $\ell < 3$.}
            \For{$k = 2^{\ell - alpha - 2} - 1$ \textbf{down to} $1$}
                \State{$k_0 \leftarrow 2^{\alpha + 1} k$}
                \For{$j = 0$ \textbf{to} $2^\alpha - 1$}
                    \State{$a[k_0 + j + 2^\alpha]
                        \leftarrow a[k_0 + j + 2^\alpha]
                        + a[k_0 + j - 2^\alpha]$}
                \EndFor
            \EndFor
            \For{$j = 0$ \textbf{to} $2^\alpha - 1$}
                \State{$a[j + 2^\alpha] \leftarrow 2 a[j + 2^\alpha]$}
            \EndFor
        \EndFor

        \For{$j = 0$ \textbf{to} $n/4 - 1$}
            \State{$(x, y) \leftarrow (a[j], a[j + n/4])$\Comment{If $a$ is
                integral, values are converted to floating-point here.}}
            \State{$z \leftarrow y \sqrt{2}$\Comment{$\sqrt{2} = 2R[1]$}}
            \State{$(a[j], a[j + n/4]) \leftarrow (x + z, x - z)$}
        \EndFor

        \For{$\alpha = 1$ \textbf{to} $\ell - 2$}
                \Comment{Loop skipped if $\ell < 3$.}
            \State{$t \leftarrow 2^{\ell - \alpha - 2}$}
            \State{$m \leftarrow 2^\alpha$}
            \For{$j = 0$ \textbf{to} $m - 1$}
                \State{$s \leftarrow R[m + j]$\Comment{$R[m + j]
                    = 1/(2\cos((4\text{\textsf{rev}}(\alpha, j)
                        + 1)\pi/2^{\alpha + 2}))$}}
                \For{$k = 0$ to $t - 1$}
                    \State{$k_0 \leftarrow 2^{\ell - \alpha - 1} j + k$}
                    \State{$k_1 \leftarrow k_0 + t$}
                    \State{$(x, y) \leftarrow (a[k_0], a[k_1])$}
                    \State{$z \leftarrow y s$}
                    \State{$(a[k_0], a[k_1]) \leftarrow (x + z, x - z)$}
                \EndFor
            \EndFor
        \EndFor
    \end{algorithmic}
\end{algorithm}

In \textsf{FFT-selfadj-new}, the first loop (lines 1-8) implements the
splits and multiplications by $1 + X$, and the third loop (lines 13-24)
performs the recombinations; however, the first loop lacks its last
iteration (last value of $\alpha$ is $\ell - 3$, not $\ell - 2$), and
the third loop lacks its first iteration (first value of $\alpha$ is 1,
not 0): these missing steps have been merged into the second loop (lines
9-12). Indeed, the last (missing) iteration of the first loop would just
have multiplied half of the value by $2$, while the first (missing)
iteration of the third loop would have multiplied the same values by
$R[1] = \sqrt{2}/2$; the two operations can be done with a single
multiplication by $\sqrt{2}$. The main reason for this merge is to
highlight an important characteristic: if the source polynomial has
integer coefficients (i.e. $a \in \bZ[X]/\Phi_n$), then all operations
in the first loop can be performed over plain integers, which is vastly
faster than using floating-point on our target hardware. As we will see,
this is the case in Falcon.

The first loop has $\ell - 2$ outer iterations, each involving $n/4$
additions, though some of these additions are really doublings; if the
source polynomial is integral, then all these operations are on
integers. The second loop uses $n/4$ addition-subtractions and $n/4$
multiplications. The third loop has $\ell - 2$ outer iterations, each
using $n/4$ addition-subtractions and $n/4$ multiplications. The total
cost of \textsf{FFT-selfadj-new} is then $(n/4)((\ell - 2)\text{\textsc{a}}
+ (\ell - 1)\text{\textsc{as}} + (\ell - 1)\text{\textsc{m}})$; if the
input is integral, the $(\ell - 2)\text{\textsc{a}}$ term disappears
(it is replaced with integer additions). This cost is less than half
the cost of \textsf{FFT-selfadj-classic}.

Algorithm \textsf{invFFT-selfadj-new}
(algorithm~\ref{alg:ifft-selfadj-new}) is a straightforward reversal of
the operations of \textsf{FFT-selfadj-new}. It uses the constants $R'[]$
which are such that $R'[j] = 1/(2R[j])$ for all $j$ (the size of $R'$ is
4096 bytes when using 64-bit floating-point values and supporting
degrees up to $n = 1024$). When the polynomial is known to be integral,
floating-point values can be converted back to integers (with
round-to-nearest) in the middle loop, and the third loop (the reverse of
the first loop in \textsf{FFT-selfadj-new}) uses only operations on
integers.

\begin{algorithm}
    \caption{\ \ \textsf{invFFT-selfadj-new}\label{alg:ifft-selfadj-new}}
    \begin{algorithmic}[1]
        \Require{FFT representation of $a$ in bit-reversal order (half-size)}
        \Ensure{$a[0]$ to $a[n/2 - 1]$ for $a\in\bR[X]/\Phi_n$ such
            that $n = 2^\ell \geq 4$ and $\adj{a} = a$}
        \For{$\alpha = \ell - 2$ \textbf{down to} $1$}
                \Comment{Loop skipped if $\ell < 3$.}
            \State{$t \leftarrow 2^{\ell - \alpha - 2}$}
            \State{$m \leftarrow 2^\alpha$}
            \For{$j = 0$ \textbf{to} $m - 1$}
                \State{$s \leftarrow R'[m + j]$\Comment{$R'[m + j]
                    = \cos((4\text{\textsf{rev}}(\alpha, j)
                        + 1)\pi/2^{\alpha + 2}) = 1/(2R[m + j])$}}
                \For{$k = 0$ to $t - 1$}
                    \State{$k_0 \leftarrow 2^{\ell - \alpha - 1} j + k$}
                    \State{$k_1 \leftarrow k_0 + t$}
                    \State{$(x, y) \leftarrow (a[k_0], a[k_1])$}
                    \State{$(u, v) \leftarrow (x + y, x - y)$}
                    \State{$(a[k_0], a[k_1]) \leftarrow (u/2, vs)$}
                \EndFor
            \EndFor
        \EndFor

        \For{$j = 0$ \textbf{to} $n/4 - 1$}
            \State{$(x, y) \leftarrow (a[j], a[j + n/4])$}
            \State{$(u, v) \leftarrow (x + y, x - y)$}
            \State{$(a[j], a[j + n/4]) \leftarrow (u/2, v(1/\sqrt{8})$
                \Comment{If $a$ is integral, values are converted to
                integers here.}}
        \EndFor

        \For{$\alpha = \ell - 3$ \textbf{down to} $0$}
                \Comment{Loop skipped if $\ell < 3$.}
            \For{$j = 0$ \textbf{to} $2^\alpha - 1$}
                \State{$a[j + 2^\alpha] \leftarrow a[j + 2^\alpha]/2$}
            \EndFor
            \For{$k = 1$ \textbf{to} $2^{\ell - alpha - 2} - 1$}
                \State{$k_0 \leftarrow 2^{\alpha + 1} k$}
                \For{$j = 0$ \textbf{to} $2^\alpha - 1$}
                    \State{$a[k_0 + j + 2^\alpha]
                        \leftarrow a[k_0 + j + 2^\alpha]
                        - a[k_0 + j - 2^\alpha]$}
                \EndFor
            \EndFor
        \EndFor
    \end{algorithmic}
\end{algorithm}

\subsection{Improved FFT}

Algorithm \textsf{FFT-selfadj-new} is an improved version of the
computation of the FFT for self-adjoint polynomials. Here, we show how
to also improve the general case, not just self-adjoint polynomials. The
main technique is to decompose the source polynomial into a sum of two
polynomials which can both be represented by self-adjoint polynomials:
given $c\in\bR[X]/\Phi_n$, we find self-adjoint polynomials $a$ and $b$
such that $c = a + X^{n/2}b$. We can then apply \textsf{FFT-selfadj-new}
on both $a$ and $b$. Using $\hat a$, $\hat b$ and $\hat c$ to denote
the FFT representations of $a$, $b$ and $c$, we have:
\begin{equation*}
    {\hat c}(\zeta) = {\hat a}(\zeta) + \zeta^{n/2} {\hat b}(\zeta)
\end{equation*}

The FFT representation returned by \textsf{FFT-classic} keeps only
$a(\zeta_{n,2j})$ for $j = 0$ to $n/2 - 1$ (specifically, it returns
only the first $n/2$ values of $a(\zeta_{n,\text{\textsf{rev}}(\ell,j')})$
for increasing values of $j'$, and the bit-reversal implies that all
these $\text{\textsf{rev}}(\ell,j')$ are even integers). We have
$\zeta_{n,2j}^{n/2} = e^{i(4j+1)\pi/2} = i$ for all $j = 0$ to $n/2$,
thus ${\hat c}(\zeta) = {\hat a}(\zeta) + i {\hat b}(\zeta)$ for all the
$\zeta$ we are interested in. Since $a(\zeta)$ and $b(\zeta)$ are all
real ($a$ and $b$ are self-adjoint), this implies that the FFT
representations of $a$ and $b$ are really the real and imaginary parts
of the elements of the FFT representation of $c$. In our chosen memory
layout with real and imaginary parts separated by $n/2$ elements, ${\hat
c}$ is then simply the concatenation of ${\hat a}$ and ${\hat b}$.

Note that $X^{n/2}b$ is neg-adjoint; we can thus easily obtain $a$ and $b$:
\begin{eqnarray*}
    a &=& (\adj{c} + c)/2 \\
    b &=& (\adj{c} - c)X^{n/2}/2
\end{eqnarray*}
Algorithm \textsf{FFT-new} (algorithm~\ref{alg:fft-new}) implements
this method; the reverse operation is performed by \textsf{invFFT-new}
(algorithm~\ref{alg:ifft-new}).

\begin{algorithm}
    \caption{\ \ \textsf{FFT-new}\label{alg:fft-new}}
    \begin{algorithmic}[1]
        \Require{$c\in\bR[X]/\Phi_n$ such that $n = 2^\ell \geq 4$}
        \Ensure{FFT representation of $c$ in bit-reversal order}
        \State{$c[0] \leftarrow 2c[0]$}
        \State{$c[n/2] \leftarrow 2c[n/2]$}
        \For{$j = 1$ \textbf{to} $n/2 - 1$}
            \State{$(x, y) \leftarrow (c[j], c[n - j])$}
            \State{$(c[j], c[n - j]) \leftarrow (x - y, x + y)$
                \Comment{$c[0..n/2-1] = 2a$, and $c[n/2..n-1] = 2b$.}}
        \EndFor
        \State{$\text{\textsf{FFT-selfadj-new}}(c[0..n/2-1])$}
        \State{$\text{\textsf{FFT-selfadj-new}}(c[n/2..n-1])$}
        \For{$j = 0$ to $n - 1$}
            \State{$c[j] \leftarrow c[j]/2$}
        \EndFor
    \end{algorithmic}
\end{algorithm}

\begin{algorithm}
    \caption{\ \ \textsf{invFFT-new}\label{alg:ifft-new}}
    \begin{algorithmic}[1]
        \Require{FFT representation of $c$ in bit-reversal order}
        \Ensure{$c\in\bR[X]/\Phi_n$ such that $n = 2^\ell \geq 4$}
        \For{$j = 0$ to $n - 1$}
            \State{$c[j] \leftarrow 2c[j]$}
        \EndFor
        \State{$\text{\textsf{invFFT-selfadj-new}}(c[0..n/2-1])$}
        \State{$\text{\textsf{invFFT-selfadj-new}}(c[n/2..n-1])$}
        \State{$c[0] \leftarrow c[0]/2$}
        \State{$c[n/2] \leftarrow c[n/2]/2$}
        \For{$j = 1$ \textbf{to} $n/2 - 1$}
            \State{$(x, y) \leftarrow (c[j], c[n - j])$}
            \State{$(c[j], c[n - j]) \leftarrow ((y + x)/2, (y - x)/2)$}
        \EndFor
    \end{algorithmic}
\end{algorithm}

In these implementations, we apply a factor of 2 to $a$
and $b$, so that the initial loop decomposes $c$ into $2a$ and $2b$:
this is done so that when $c$ is an integral polynomial (as is the
case in Falcon) then the decomposition can be performed with only
integer operations. The corrective factor $1/2$ is applied on the
FFT representation (halving is inexpensive on floating-point values).

Total cost of \textsf{FFT-new} is twice that of \textsf{FFT-selfadj-new},
plus an extra $n/2$ addition-subtractions when the the source polynomial
is not integral, hence $(n/2)((\ell - 2)\text{\textsc{a}}
+ \ell\text{\textsc{AS}} + (\ell - 1)\text{\textsc{m}})$ in the general
case, which is less than half the cost of \textsf{FFT-classic}. The
inverse algorithm \textsf{invFFT-new} has the same cost.

\begin{remark}
There are other possible decompositions of a polynomial $c$ into two
self-adjoint polynomials; for instance, one may find self-adjoint
polynomials $a$ and $b$ such that $c = a + b/(1 + X)$, i.e. $c$ is the
sum of a self-adjoint and an $X$-adjoint polynomials. We here favour the
$c = a + X^{n/2} b$ decomposition because it yields the same FFT
representation as \textsf{FFT-classic}, thus keeping other operations in
Falcon unchanged.
\end{remark}

\subsection{Input Range and Benchmarks}\label{sec:fft-new-range}

We focus here on \textsf{FFT-new}, with an integral input polynomial
whose coefficients are represented over 32-bit integers. The
computations which are performed over plain integers are correct only as
long as they do not incur an overflow. In the opening loop of
\textsf{FFT-new} (algorithm~\ref{alg:fft-new}, steps 1-5), each value
$c[j]$ is replaced with the sum or the difference of two coefficients of
the input. The resulting values are then passed to
\textsf{FFT-selfadj-new} (algorithm~\ref{alg:fft-selfadj-new}, steps
1-7) which runs $\ell - 2$ iterations of a loop; at each iteration, some
of the values are replaced with the sum of two values produced by the
previous iteration. In total, this means that $|c[j]| \leq 2^{\ell-1}
\norminf{c}$ for any $j$, when reaching step 8 of
\textsf{FFT-selfadj-new}; afterwards, values are converted to
floating-point and no longer limited to the 32-bit integer range.
Correspondingly, for degree $n \leq 1024$, if $\norminf{c} < 2^{22}$,
then no overflow may occur in the entire \textsf{FFT-new}. All uses of
\textsf{FFT-new} (and \textsf{FFT-selfadj-new}) in this paper will be
shown to fulfill that condition.

We implemented and benchmarked these algorithms on an Arm Cortex M4
microcontroller (code is in C but the addition, addition-subtraction and
multiplications are hand-optimized in assembly). For an integral
polynomial held in 32-bit signed integers, \textsf{FFT-classic} runs in
611608 cycles (for degree $n = 512$), or 1369401 cycles (for $n =
1024$). With \textsf{FFT-new}, these figures drop to 385940 and 855782
cycles, respectively. We can see that the extra integer-only operations
have a non-negligible cost, but the performance gain is nonetheless
substantial.

\subsection{Annex: Extension to NTT}

The \emph{Number Theoretic Transform} (NTT) is analogous to the FFT when
working modulo a prime integer $q$. To apply the NTT on polynomials in
$\bZ_q[X]/\Phi_n$, we normally want $\Phi_n = X^n + 1$ to split over
$Z_q$, i.e. that $q = 1 \bmod 2n$. This is a slightly restrictive
condition. For instance, Falcon uses $q = 12289$ because it needs to
work up to degree $n = 1024$, and 12289 is the smallest prime that is
equal to 1 modulo 2048. It is possible to some extent to use a
\emph{partial} NTT if the modulus is not adequate for the full degree;
for instance, ML-KEM\cite{Fips203} uses $q = 3329$, which is good for NTT
up to degree 128, but applies it to polynomials modulo $X^{256} + 1$.
The method is simply to split the input $a$ into $a_{\text{even}}$ and
$a_{\text{odd}}$ and apply the NTT on both; however, multiplications of
polynomials in such a partial NTT representation become more
complicated:
\begin{equation*}
    (ab)(\zeta) =
        (a_{\text{even}} b_{\text{even}}
            + X a_{\text{odd}} b_{\text{odd}})(\zeta^2)
        + \zeta (a_{\text{even}} b_{\text{odd}}
            + a_{\text{odd}} b_{\text{even}})(\zeta^2)
\end{equation*}
i.e. more operations in $\bZ_q$ are needed compared to a simple
coefficient-wise multiplication. Generally speaking, if we want to work
with polynomials of degree $n = 2^\ell$ but the modulus allows the plain
NTT only up to $2^{\ell - \gamma}$ then multiplications of two
polynomials in partial NTT representation will need $O(2^\gamma n)$
multiplications in $\bZ_q$. We thus prefer $\gamma$ to remain small. For
degree $n = 1024$, there are not many good choices of primes lower than
12289: 7681, 10753 and 11777 allow up to degree 256 (hence $2^\gamma =
4$), while 257, 769, 3329, 7937 and 9473 work natively only up to degree
128, and thus imply $2^\gamma = 8$ when $n = 1024$.

A different way is to consider $q = -1 \bmod n$. With such a modulus,
$-1$ is not a square in $\bZ_q$, and we can extend that field to
$\bF_{q^2}$ by adjoining a formal element $i$ such that $i^2 = -1$; this
is the same construction as complex numbers with regard to reals. Note
that if $q = -1 \bmod n$ then there is some integer $k$ such that $q =
kn - 1$, which implies that $q^2 - 1 = (q - 1)(q + 1) = (kn - 2)kn$.
With $n = 2^\ell$ being even, we get that $q^2 = 1 \bmod 2n$: polynomial
$\Phi_n$ splits over $\bF_{q^2}$, and we can apply the FFT, including
the improved \textsf{FFT-new} and \textsf{invFFT-new} algorithms. All
operations on complex numbers, in particular conjugation, directly map
to $\bF_{q^2}$. We only need to arbitrarily choose a root $\delta$ of
$\Phi_n = X^n + 1$ in $\bF_{q^2}$, then use the roots $\zeta_{n,j} =
\delta^{2j+1}$ for $j = 0$ to $n - 1$. In total, we get an FFT and
inverse FFT with good efficiency (comparable to a plain full-degree
NTT), and multiplications of two polynomials in FFT representation (in
$\bF_{q^2}$) require $2n$ multiplications in $\bF_q$ (the equivalent of
$2^\gamma = 2$ in terms of efficiency). Available moduli less than 12289
and that support degree $n = 1024$ for this FFT in $\bF_{q^2}$ are 5519,
6143 and 8191; this last value, in particular, is a Mersenne prime, for
which modular reduction is easy, making it attractive from an
implementation point of view. If we accept to use a partial FFT, which
increases the cost of polynomial multiplication, we can use more moduli,
e.g. 1279 (good up to degree 256 natively) or 127 (good up to degree
128).

Falcon itself is specified only with $q = 12289$, and the plain NTT is
the way to go for that modulus. But the FFT in $\bF_{q^2}$, and its
improved implementation \textsf{FFT-new}, can be useful for designing
new schemes (or new parameter sets for Falcon) since they offer
reasonably fast operations over additional moduli that help strike a
better trade-off between object sizes (e.g. encoded public key) and
security.

% ----------------------------------------------------------------------

\newpage
\section{Falcon Tree Optimization}\label{sec:tree}

While we primarily target RAM optimizations, we are also happy to reduce
runtime cost. Floating-point values are both large (8 bytes each) and
expensive, thus we would like to have fewer of them. The \emph{Falcon
tree} is a structure that depends only on the private key; a
speed-optimized implementation for an application where it is expected
that many messages will be signed with the same key (e.g. a TLS server)
will usually precompute the tree when the private key is loaded; in that
case, the time taken to build this tree does not matter much (since the
cost is amortized over many signatures), but the storage size could be
important. On the other hand, a RAM-optimized implementation, in
particular in the context of small embedded systems (for which RAM is
usually a very scarce resource), will usually resort to dynamically
generating the Falcon tree along with each signature, keeping in RAM at
any point only the relevant parts of the tree (i.e. the path from the
root to the particular leaf which is currently used). In this section,
we present some optimization techniques that help in both cases.

\subsection{\emph{1/X}-Adjoint Polynomials}

The \textsf{FalconTreeTop} algorithm (algorithm~\ref{alg:FalconTreeTop})
generates the root node of the Falcon tree, then uses
\textsf{FalconTreeInner} (algorithm~\ref{alg:FalconTreeInner}) to build
the left and right sub-trees. By construction, the parameter to
\textsf{FalconTreeInner} is a self-adjoint polynomial $a$, which the
algorithm splits: $(b, c) = \text{\textsf{split}}(a)$. As was previously
noticed (see remark~\ref{remark:split-selfadj}), splitting a
self-adjoint polynomial yields a self-adjoint and an $X$-adjoint
polynomials, i.e. $\adj{b} = b$ and $\adj{c} = Xc$. Polynomial $\adj{c}$
is itself $1/X$-adjoint, since $\adj{\adj{c}} = c = (1/X)Xc =
(1/X)\adj{c}$. The root of the sub-tree returned by
\textsf{FalconTreeInner} is the value $L = \adj{c}/b$, which is then an
$1/X$-adjoint polynomial.

Previous implementations of Falcon stored $L$ as a generic polynomial in
FFT representation, using $m/2$ floating-point elements ($4m$ bytes) for
an input polynomial $a \in \bR[X]/\Phi_m$ (the \textsf{split} operation
then implies that $b$ and $c$, and thus $L$, have degree $m/2$).
However, we can instead store $L/(1 + X)$, which is self-adjoint and
requires only $m/4$ floating-point elements ($2m$ bytes). Algorithm
\textsf{LDL-new} (algorithm~\ref{alg:ldl-new}) takes as input $a$
(self-adjoint, FFT representation), splits it into $b$ (i.e.
$a_{\text{even}}$) and $c$ (i.e. $a_{\text{odd}}$), then returns
polynomials $d$ and $\lambda$ (both self-adjoint and in FFT
representation) such that:
\begin{eqnarray*}
    d &=& b - c\adj{c}/b \\
    \lambda &=& L/(1+X)
\end{eqnarray*}
for $L = \adj{c}/b$. This is mostly what \textsf{FalconTreeInner}
computes, except for the specific representation of $L$ as the
self-adjoint polynomial $\lambda$.

\begin{algorithm}
    \caption{\ \ \textsf{LDL-new}\label{alg:ldl-new}}
    \begin{algorithmic}[1]
        \Require{$a \in \bR[X]/\Phi_m$ such that $a = \adj{a}$ (FFT,
            half-size)}
        \Ensure{$b$, $d = b - c\adj{c}/b$ and $\lambda = \adj{c}/(b(1+X))$ for
            $(b, c) = \text{\textsf{split}}(a)$ (FFT, degree $m/2$, half-size)}
        \For{$j = 0$ \textbf{to} $m/4 - 1$}
            \State{$(x, y) \leftarrow (a[2j], a[2j + 1])$}
            \State{$u \leftarrow (x + y)/2$}
            \State{$v \leftarrow (x - y)R'[m/4 + j]$\Comment{Same
                constants $R'[j]$ as in \textsf{invFFT-selfadj-new}.}}
            \State{$\lambda[j] \leftarrow v / (u C[m/4 + j])$}
            \State{$d[j] \leftarrow u - v \lambda[j]$}
        \EndFor
    \end{algorithmic}
\end{algorithm}

In \textsf{LDL-new}, steps 2-4 represent $(b, c) =
\text{\textsf{split}}(a)$, with $c$ being represented by $c' = (1+X)c$;
then, variables $u$ and $v$ are equal to $b[j]$ and $c'[j]$,
respectively ($b$ and $c'$ are self-adjoint). Then, we want:
\begin{equation*}
    \lambda = \frac{\adj{c}}{(1+X)b} = \frac{c'}{\adj{(1+X)}(1+X)b}
        = \left(\frac{X}{(1+X)^2}\right)\frac{c'}{b}
\end{equation*}
Hence, the constant $C[m/4 + j]$ is equal to $(1+X)^2/X$ evaluated over
the roots $\zeta$ of $\Phi_{m/2}$ (equivalenty, the roots $\zeta^2$ of
$\Phi_m$) in the correct bit-reversal order, namely:
\begin{equation*}
    C[2^{\alpha - 1} + \text{\textsf{rev}}(\alpha - 1, j)]
        = 2 + 2\cos((4j + 1)\pi/2^\alpha)
\end{equation*}
for $j = 0$ to $2^{\alpha - 1}$. To support Falcon up to degree
$n = 2^\ell$, these constants must be
precomputed for $\alpha$ up to $\ell - 1$ (since the largest $d$ and
$\lambda$ are for degree $n/2$), for a total size of 4096 bytes for
$n = 1024$ and using 64-bit floating-point values.

When using the Falcon tree in \textsf{ffSampling}, the tree node value $L$
must be multiplied with a generic polynomial, both being in FFT
representation (algorithm~\ref{alg:ffSampling}, step 10). Since we do not
have $L$, but $\lambda = L/(1+X)$, this multiplication is modified;
algorithm \textsf{mul-iXadj} (algorithm~\ref{alg:mul-iXadj}) performs
this operation.

\begin{algorithm}
    \caption{\ \ \textsf{mul-iXadj}\label{alg:mul-iXadj}}
    \begin{algorithmic}[1]
        \Require{$a \in \bR[X]/\Phi_m$ (FFT), $\lambda \in \bR[X]/\Phi_m$
            such that $\adj{\lambda} = \lambda$ (FFT, half-size)}
        \Ensure{$d = (1 + X)\lambda a$ (FFT)}
        \For{$j = 0$ \textbf{to} $n/2 - 1$}
            \State{$(x_r, x_i) \leftarrow (a[j], a[j + n/2])$}
            \State{$y \leftarrow \lambda[j]$}
            \State{$\theta = y C[n/2 + j]/2$\Comment{Same constants $C[j]$ as
                in \textsf{LDL-new}}}
            \State{$\delta = D[n/2 + j]$}
            \State{$d[j] \leftarrow \theta(x_r + x_i \delta)$}
            \State{$d[j + n/2] \leftarrow \theta(x_i - x_r \delta)$}
        \EndFor
    \end{algorithmic}
\end{algorithm}

This uses the constants $D[j]$ which are equal to
$i(\zeta - 1)/(\zeta + 1)$ for the roots $\zeta$ of $\Phi_m$; in
the appropriate bit-reversal order:
\begin{eqnarray*}
    D[2^{\alpha - 1} + \text{\textsf{rev}}(\alpha - 1, j)]
        = \frac{-\sin((4j + 1)\pi/2^\alpha)}{1 + \cos((4j + 1)\pi/2^\alpha)}
\end{eqnarray*}
Again, the $D[]$ array size is 4096 bytes when supporting Falcon up to
$n = 1024$ and using 64-bit floating-point values.

\subsection{Symplecticity}

Symplecticity is a geometric property fulfilled by Falcon
trees\cite{SunZhoZhaTaoQiaMin2022}. It provides some useful relations
between the left and right sub-trees of the root nodes; namely, the
right sub-tree can be obtained from the left sub-tree by doing the
following:
\begin{itemize}
    \item Negate all values ($L$) of intermediate nodes.

    \item Exchange left and right children of each node.

    \item Replace each leaf value $x$ with $q^2/x$.
\end{itemize}
If the Falcon tree is precomputed and stored in RAM, then these properties
reduce both the storage size and the computational cost of building the
tree; values of inner nodes of the right sub-tree need not be stored since
they can be obtained by using the mirror nodes in the left sub-tree, with
a very inexpensive negation, and the $n/2$ leaves of the right sub-tree
are computed directly from the $n/2$ leaves of the left sub-tree.

The storage cost of a full Falcon tree is then:
\begin{itemize}

    \item Top-level
    $\text{\textsf{T.value}} = (F\adj{f} + G\adj{g})/(f\adj{f} + g\adj{g})$:
    $n$ floating-point values.

    \item Left sub-tree: $\ell - 1$ layers, each containing
    $2^{\ell - 1 - \alpha}$ values $\lambda \in\bR[X]/\Phi_{2^\alpha}$
    which are self-adjoint, for a total of $n/4$ floating-point values,
    and $n/2$ leaves.

    \item Right sub-tree: $n/2$ leaves (the inner nodes are omitted).
\end{itemize}
This means a total of $n(\ell + 7)/4$ floating-point values. If the
values are 64-bit objects, then this means 16384 bytes (for $n = 512$)
or 34816 bytes (for $n = 1024$).

In our implementation, we try to achieve low RAM usage, and thus we do
not precompute the Falcon tree; instead, we compute it dynamically,
keeping only a single path in RAM at a time, which implies that we have
to recompute the $L$ values in the right sub-tree. However, we can still
leverage a consequence of symplecticity: the source value for the right
sub-tree is the polynomial denoted $d$ in \textsf{FalconTreeTop}
(algorithm~\ref{alg:FalconTreeTop}, step 8); this value is such that:
\begin{eqnarray*}
    (f\adj{f} + g\adj{g}) d
      &=& (F\adj{F} + G\adj{G})(f\adj{f} + g\adj{g})
        - (F\adj{f} + G\adj{g})(\adj{F}f + \adj{G}{g}) \\
      &=& gF\adj{(gF)} + fG\adj{(fG)} - gF\adj{(fG)} - fG\adj{(gF)} \\
      &=& (gF - fG)\adj{(gF - fG)} \\
      &=& q^2
\end{eqnarray*}
by application of the NTRU equation. We can thus obtain the source
value for the right sub-tree as $d = q^2/(f\adj{f} + g\adj{g}) = q^2/a$,
with $a$ being the source value for the left sub-tree. Since $a$ is
in FFT representation and is self-adjoint, we get $d$ with $n/2$
divisions on floating-point values.

An alternative way would be to specialize the left and right
\textsf{FalconTreeInner} so that the dynamic reconstruction of the right
sub-tree runs, in fact, the left sub-tree building and applies the
symplectic negations and swaps, and the final divisions on the leaves.
However, that would complicate the implementation (inner
\textsf{ffSampling} calls would have to know whether they are exploring
the left or right sub-tree); computing $d = q^2/a$ is simpler. Note that
both methods involve the same number ($n/2$) of floating-point
divisions.

\subsection{Top-Level Polynomials}\label{toplevel-range}

\textsf{FalconTreeTop} (algorithm~\ref{alg:FalconTreeTop}) computes the
top-level tree value $L$, and invokes \textsf{FalconTreeInner} on two
polynomials equal to $a = f\adj{f} + g\adj{g}$ and, as was pointed out
above, $d = q^2/a$. The top-level tree value is $L = b/a$. We therefore
need to compute the two polynomials $a$ and $b$. Note that both
polynomials are integral; we can thus evaluate them over integers, and
more precisely we compute $a$ and $b$ modulo $q = 12289$ and modulo $q'
= 18433$, then assemble the resulting coefficients with the CRT to yield
the result. This process works as long as the integer coefficients can
be shown to be less than $qq'/2$ in absolute value. In this subsection,
we establish the maximum range of coefficients of $a$ and $b$. We rely
on two properties of valid Falcon private keys:
\begin{eqnarray*}
    \norminf{(F,G)} &\leq& 2^7 \\
    \norm{(g, -f)} &\leq& 1.17\sqrt{q}
\end{eqnarray*}
The second property can also be expressed as follows:
\begin{equation*}
    \sum_{j=0}^{n-1} f[j]^2 + \sum_{j=0}^{n-1} g[j]^2 \leq 16822
\end{equation*}

Let's first consider $b = F\adj{f} + G\adj{g}$. For the purpose of range
analysis, we can consider that all coefficients of $F$ and $G$ are equal
to the upper bound of $\norminf{(F, G)}$, which is $2^7$; hence:
\begin{eqnarray*}
    \norminf{F\adj{f} + G\adj{g}} \leq 2^7 \sum_{j=0}^{n-1} (|f[j]| + |g[j]|)
\end{eqnarray*}
We thus want to upper bound the sum of the (absolute values) of
coefficients of $f$ and $g$, under the condition that
$\norm{(g,-f)}^2 \leq 16822$. The following lemma will be useful:

\begin{lemma}\label{lemma:max-sum}
Let $m \in\bN$. Let $S_m = \{ (x, y)\in\bN^2 \,|\, x^2 + y^2 \leq m \}$.
Then the maximum value of $x + y$ for $(x, y)\in S_m$ is achieved for some
integers $x$ and $y$ such that $|x - y|\leq 1$.
\end{lemma}
\begin{proof}
Let $(x, y)\in S_m$ that maximizes $x + y$. If $y \geq x + 2$, then we
define $(x', y') = (x + 1, y - 1)$. It holds that
$x'^2 + y'^2 = x^2 + y^2 + 2(x - y) + 2 < x^2 + y^2$,
therefore $(x', y')\in S_m$. However, $x' + y' = x + y$, thus $(x', y')$
is also a maximizing pair. Applying this process repeatedly, we can find
a pair $(x'', y'')\in S_m$ such that $|x'' - y''| \leq 1$ that
also maximizes the sum $x'' + y''$.\qed
\end{proof}

For any pair $(f,g)$ of polynomials that fulfills the condition
$\norm{(g,-f)}^2 \leq 16822$, we can consider the smallest and largest of
the $2n$ coefficients (in absolute value); if the difference between the
two is 2 or more, then we can apply lemma~\ref{lemma:max-sum} and obtain
another pair $(f',g')$ which is almost identical to $(f,g)$ except that
the two aforementioned coefficients have been replaced with values $(x,
y)$ such that $|x - y|\leq 1$, without changing the total sum
$\sum_{j=0}^{n-1} (|f[j]| + |g[j]|)$, and still complying with the
condition on $\norm{(g,-f)}$. Doing so repeatedly shows that under that
condition, $\sum_{j=0}^{n-1} (|f[j]| + |g[j]|)$ can be maximized when all
coefficients of $f$ and $g$ are equal to either $x$ or $x+1$ for some
integer $x$.

We can thus write that:
\begin{equation*}
    \norminf{F\adj{f} + G\adj{g}} \leq 2^7 ((2n - m) x + m (x + 1))
\end{equation*}
for some integers $m\in [0, 2n-1]$ and $x$ which are such that $(2n - m)
x^2 + m (x + 1)^2 \leq 16822$. Note that $(2n - m) x + m(x + 1) = 2nx +
m$, which is maximized for any given $x$ when $m$ is maximized, under
the constraint that $(2n - m) x^2 + m (x + 1)^2 \leq 16822$, which means
that we need $m \leq (16822 - 2nx^2)/(2x + 1)$. Trying increasing values
of $x$ until there is no matching $m$, we can easily enumerate the
largest possible value for $(2n - m) x + m (x + 1)$, which is 5822 for
$n = 1024$, and 4144 for $n = 512$ (this upper bound is smaller for lower
degrees). Correspondibly, this demonstrates that:
\begin{equation*}
    \norminf{F\adj{f} + G\adj{g}} \leq 5822\cdot 2^7 = 745216
\end{equation*}
which is both lower than $qq'/2$ (hence we can find the correct value
of $F\adj{f} + G\adj{g}$ by working only modulo $q$ and modulo $q'$),
and lower than $2^{22}$, which implies that $F\adj{f} + G\adj{g}$ is
an adequate input to algorithm \textsf{FFT-new} with 32-bit integers
(see section~\ref{sec:fft-new-range}).

For $a = f\adj{f} + g\adj{g}$, we can prove a tighter range, using
the following lemma:
\begin{lemma}\label{lemma:ffa-norm}
For any polynomial $f \in \bR[X]/\Phi_n$, $\norminf{f\adj{f}} = \norm{f}^2$.
\end{lemma}
\begin{proof}
Let $k\in [1, n-1]$. A straightforward computation of $f\adj{f}$ shows
that $f\adj{f}[k]$ is the sum of values $\pm f[j] \adj{f}[k - j]$, with
the convention that $f[j - n] = -f[j]$ (which corresponds to reduction
modulo $X^n + 1$). By definition of the adjoint,
$|\adj{f}[k - j]| = |f[j - k]|$. By triangular inequality, we have:
\begin{equation*}
    |f\adj{f}[k]| \leq \sum_{j=0}^{n-1} |f[j]| |f[j - k]|
\end{equation*}
For any reals $x$ and $y$, we have $x^2 + y^2 - 2xy = (x - y)^2 \geq 0$,
thus $xy \leq (x^2 + y^2)/2$. Applying this rule to the sum above, we
get:
\begin{equation*}
    |f\adj{f}[k]| \leq \sum_{j=0}^{n-1} \frac{f[j]^2 + f[j - k]^2}{2}
\end{equation*}
Each $f[j]^2$ appears exactly twice in the right-hand side sum, hence:
\begin{equation*}
    |f\adj{f}[k]| \leq \sum_{j=0}^{n-1} f[j]^2 = \norm{f}^2
\end{equation*}
for all $k\in [1, n-1]$. At index 0, the inequality becomes an equality:
\begin{equation*}
    f\adj{f}[0] = f[0]^2 - \sum_{j=1}^{n-1} f[j] \adj{f}[n - j]
        = f[0]^2 + \sum_{j=1}^{n-1} f[j]^2 = \norm{f}^2
\end{equation*}
which completes the proof.\qed
\end{proof}

From lemma~\ref{lemma:ffa-norm}, we get that:
\begin{equation*}
    \norminf{f\adj{f} + g\adj{g}} \leq \norminf{f\adj{f}} + \norminf{g\adj{g}}
        = \norm{f}^2 + \norm{g}^2 = \norm{(g,-f)}^2
\end{equation*}
Thus, all coefficients of $f\adj{f} + g\adj{g}$ are at most 16822 (in
absolute value).

If we had $q' > 2\cdot 16822$, then we could compute $f\adj{f} +
g\adj{g}$ modulo $q'$ only; it would be sufficient. However, that would
require $q' > 33644$ (next prime appropriate for the NTT at $n = 1024$
is 40961), and there are computational advantages to sticking to a
modulus lower than $2^{15} = 32768$ (in particular, this makes it easier
to use some SIMD instructions of the Arm Cortex M4 such as \verb+smulbb+
which work only with \emph{signed} 16-bit operands). We will therefore
use $q' = 18433$ and compute $f\adj{f} + g\adj{g}$ modulo $q$ and modulo
$q'$. The reconstruction of each coefficient as a plain integer is
easier than with the CRT, though. Suppose that for a given index $j$, we
obtain the coefficient modulo $q$ as a value that normalizes as $x \in
[-q/2, +q/2]$, and similarly the same coefficient modulo $q'$ normalizes
as $y \in [-q'/2, +q'/2]$. If $x = y$, $x = y + q$ or $x = y - q$, then
$x$ agrees with $y$, which is the value of the coefficient as an
integer. Otherwise, the value $y$ must be adjusted: if $y < 0$, then the
correct value is $y + q'$, otherwise the correct value is $y - q'$.
These operations are easy and inexpensive to apply. Moreover, if we have
to compute $f\adj{f} + g\adj{g}$ twice (we will see that this is the
case in our implementation, for RAM-saving reasons), then we can
remember whether an adjustment to the value of the coefficient modulo
$q'$ was needed, and the second time we only compute the polynomial
modulo $q'$ and apply the adjustement based on the sign of each value.
The storage space for such remembering is small ($n$ bits, i.e. $n/8$
bytes); in our implementation, we store that bitfield temporarily in the
output buffer for the signature.

% ----------------------------------------------------------------------

\newpage
\section{Signature Generation Optimization}\label{sec:genopt}

\subsection{Optimization Techniques}

We here present the main techniques we use to achieve a small RAM usage,
as well as better performance on Arm Cortex M4 CPUs. The following ideas
are used:
\begin{itemize}

    \item Use \textsf{FFT-new} instead of \textsf{FFT-classic}.

    \item Perform computations on integer polynomials using modular
    integers, working modulo $q = 12289$ and $q' = 18433$.

    \item Build the Falcon tree dynamically, working only with
    self-adjoint polynomials (which use only half the storage space),
    as described in the previous section.

    \item Use only the fractional part of the input to \textsf{ffSampling}.
    Moreover, that input can be injected later in the process, at the
    leaf level in the Falcon tree rather than at the start, which avoids
    some FFTs.

    \item Record sampled integers as returned by \textsf{SamplerZ} and
    use these values to compute the signature, rather than the final
    polynomial returned by \textsf{ffSampling}. This removes the need to
    run an inverse FFT. In fact, \textsf{invFFT-new} is not used at all
    in our implementation.

\end{itemize}

\paragraph{Fractional Input.}
A well-known property of \textsf{ffSampling} is that it preserves the
integral part of its input (\cite{AlmFouLacPre2026}, lemma 9): if we
consider a Falcon tree $\text{\textsf{T}}$, a given state for the
pseudorandom generator used in \textsf{SamplerZ}, and an input
$\mathbf{t} = (t_0, t_1)$, yielding $\mathbf{z} =
\text{\textsf{ffSampling}}(\mathbf{t}, \text{\textsf{T}})$, then for any
$\mathbf{k} = (k_0, k_1) \in (\bZ[X]/\Phi_n)^2$,
$\text{\textsf{ffSampling}}(\mathbf{t} + \mathbf{k}, \text{\textsf{T}})$
should return $\mathbf{z} + \mathbf{k}$ with the same $\mathbf{z}$ if
the same pseudorandom generator state is used. We can thus split the
input into an integral and a fractional parts, perform the sampling on
the fractional part only, then add back the integral part afterwards.
In Falcon signature generation, the input vector is $(t_0, t_1) =
(-cF/q, cf/q)$; we can split $cf/q = (cf \bmod q)/q + (cf - (cf \bmod
q))/q$, the second part being integral. We will thus run
\textsf{ffSampling} on $((-cF \bmod q)/q, (cf \bmod q)/q)$.

\paragraph{Late Addition.}
Algorithm \textsf{ffSampling} only repeatedly splits its input, and
splitting (in coefficient representation) is a data routing operation,
i.e. it does not actually modify the coefficient values. In
algorithm~\ref{alg:ffSampling} (step 10), the polynomial $t'_0$ is
obtained by \emph{adding} $(t_1 - z_1)L$ (produced by the previous
sampling steps) to the input $t_0$; addition commutes with splitting,
thus we can delay this addition until it is really needed, which is at
the leaf level. In other words, instead of computing $(t_0, t_1)$ and
passing it as parameter to \textsf{ffSampling}, we can run
\textsf{ffSampling} on $(0, 0)$ and add coefficients of $((-cF \bmod
q)/q, (cf \bmod q)/q)$ only at the point where \textsf{SamplerZ} is
called (algorithm~\ref{alg:ffSampling}, steps 3 and 4). We keep $-cF
\bmod q$ and $cf \bmod q$ as arrays of 16-bit signed values ($2n$ bytes
each); as is usual in splitting-based algorithms, these coefficients are
consumed in bit-reversal order. The division by $q$, and conversion to
floating-point, is only done at that point. This late addition has the
benefit of avoiding an FFT computation: the FFT representation of 0 is
0. We can also skip some of the computations in \textsf{ffSampling} by
propagating a flag that indicates that the input value is statically
known to be 0 at this point. Apart from the obvious speed benefits (not
running the FFT is always faster than running it), this also saves some
RAM space because applying the FFT on a degree-$n$ polynomial returns it
as $8n$ bytes (with 64-bit floating-point elements), while working
modulo $q$ only needs $2n$ bytes.

\paragraph{Signature Building.}
The Gaussian sampler \textsf{SamplerZ} produces integer values
(algorithm~\ref{alg:ffSampling}, steps 3 and 4); these values are
ultimately assembled into the vector $\mathbf{z} = (z_0, z_1)$ and the
overall \textsf{ffSampling} algorithm returns $\mathbf{z}$, normally in
FFT representation, which the caller subtracts from the input
$\mathbf{t}$. Using this vector entails performing an inverse FFT at
some point. However, we can also discard it entirely; instead, we can
record the coefficients of the sampled $z_0$ and $z_1$ in an alternate
buffer (for instance, the same buffer that provides the coefficients of
$-cF \bmod q$ and $cf \bmod q$ can be used to receive the sampled $z_0$
and $z_1$). The signature value is ultimately computed as:
\begin{eqnarray*}
    \mathbf{s} &=& (\mathbf{t} - \mathbf{z})\mathbf{B} \\
        &=& \mathbf{t}\mathbf{B} - \mathbf{z}\mathbf{B} \\
        &=& (c - g z_0 - G z_1, f z_0 + F z_1)
\end{eqnarray*}
These polynomials are all integral, hence we can compute the signature
by working modulo $q$ or modulo $q'$.

\paragraph{Merged Tree Building and Sampling.}
Using the techniques above, we merge the dynamic Falcon tree building
and the Fourier sampling into algorithm \textsf{ffTreeSampling}
(algorithm~\ref{alg:ffTreeSampling}). It takes as input polynomials $t$
and $a$, both from $\bR[X]/\Phi_m$; $a$ is self-adjoint and represents
the input to \textsf{FalconTreeInner}, while $t$ is the input to
\textsf{ffSampling} (before splitting).

\begin{algorithm}
    \caption{\ \ \textsf{ffTreeSampling}\label{alg:ffTreeSampling}}
    \begin{algorithmic}[1]
        \Require{$t \in \bR[X]/\Phi_m$ (FFT), $a \in \bR[X]/\Phi_m$
            (FFT, self-adjoint, half-size), I/O state $W$}
        \Ensure{$u = t - z \in \bR[X]/\Phi_m$ (FFT)}
        \If{$m = 2$}
            \State{$\sigma' \leftarrow \sigma_n/\sqrt{a[0]}$
                \Comment{Leaf normalization step.}}
            \State{$(e_0, e_1) \leftarrow \text{\textsf{next-input}}(W)$}
            \State{$z_0 \leftarrow \text{\textsf{SamplerZ}}(t[0] + e_0)$
                \Comment{Late addition of input coefficient.}}
            \State{$u[0] \leftarrow t[0] - z_0$}
            \State{$z_1 \leftarrow \text{\textsf{SamplerZ}}(t[1] + e_1)$
                \Comment{Late addition of input coefficient.}}
            \State{$u[1] \leftarrow t[1] - z_1$}
            \State{$\text{\textsf{push-sample}}(W, z_0, z_1)$
                \Comment{Sampled integers are saved.}}
        \Else
            \State{$(t_0, t_1) \leftarrow \text{\textsf{split}}(t)$}
            \State{$(b, d, \lambda) \leftarrow \text{\textsf{LDL-new}}(a)$}
            \State{$u_1 \leftarrow \text{\textsf{ffTreeSampling}}(t_1, d, W)$
                \Comment{$u_1 = t_1 - z_1$}}
            \State{$v_1 \leftarrow \text{\textsf{mul-iXadj}}(u_1, \lambda)$
                \Comment{$v_1 = (t_1 - z_1)L$}}
            \State{$t'_0 \leftarrow t_0 + v_1$
                \Comment{$v_1 = t'_0 - t_0$}}
            \State{$v_0 \leftarrow \text{\textsf{ffTreeSampling}}(t'_0, b, W)$
                \Comment{$v_0 = t'_0 - z_0$}}
            \State{$u_0 \leftarrow v_0 - v_1$
                \Comment{$u_0 = t_0 - z_0$}}
            \State{$u \leftarrow \text{\textsf{merge}}(u_0, u_1)$
                \Comment{$u = t - z$}}
        \EndIf
        \State{\Return{$u$}}
    \end{algorithmic}
\end{algorithm}

The running state $W$ encapsulates an input/output buffer $w$ for 16-bit
values, and a counter $k$. Initially, the buffer contains the
coefficient of the fractional additional input (e.g. $cf \bmod q$) and
the counter is set to $n/2 = 2^{\ell - 1}$; each invocation of
\textsf{next-input} (algorithm~\ref{alg:next-input}) returns two values
from the buffer, while each invocation of \textsf{push-sample}
(algorithm~\ref{alg:push-sample}) receives two sampled integers that
replace the two previously returned values in the buffer.

\begin{algorithm}
    \caption{\ \ \textsf{next-input}\label{alg:next-input}}
    \begin{algorithmic}[1]
        \Require{$W = (w, k)$ (buffer $w$, counter $k$)}
        \Ensure{$(e_0, e_1)$; $k$ is updated}
        \State{$k \leftarrow k - 1$}
        \State{$j \leftarrow \text{\textsf{rev}}(\ell - 1, k)$}
        \State{\Return $(w[j]/q, w[j + n/2]/q)$}
    \end{algorithmic}
\end{algorithm}

\begin{algorithm}
    \caption{\ \ \textsf{push-sample}\label{alg:push-sample}}
    \begin{algorithmic}[1]
        \Require{$W = (w, k)$ (buffer $w$, counter $k$), integers $(z_0, z_1)$}
        \Ensure{$w$ is updated}
        \State{$j \leftarrow \text{\textsf{rev}}(\ell - 1, k)$}
        \State{$w[j] \leftarrow z_0$}
        \State{$w[j + n/2] \leftarrow z_1$}
    \end{algorithmic}
\end{algorithm}

\subsection{Full Signing Procedure}

Putting these techniques together, we obtain the following signature
generation process:
\begin{enumerate}
    \item Compute the message representative $\mu$, generate the random
    seed $r$, and sample the polynomial $c\in \bZ_q[X]/\Phi_n$.

    \item Compute $a = f\adj{f} + g\adj{g}$ by working modulo $q$ and
    modulo $q'$, and rebuild its plain integer coefficients (as
    shown in section~\ref{toplevel-range}, these coefficients are
    in $[-16822, +16822]$ and thus only need 2 bytes each).

    \item Convert $a$ to FFT with \textsf{FFT-selfadj-new}, then
    compute $d = q^2/a$ (coefficient-wise floating-point division, since
    we operate on the FFT representation).

    \item Compute $w = cf \bmod q$, normalized to coefficients in
    $[-q/2, +q/2]$, and set counter $k = n/2$. $w$ and $k$ are the
    two components of state $W$.

    \item Run $\text{\textsf{ffTreeSampling}}(0, d, W)$. This yields
    $u_1 = t_1 - z_1$. Note that we provided only a fractional input
    ($(cf \bmod q)/q$, not the full $cf/q$) but the missing integral
    parts cancel out in $u_1$ (we virtually subtract the same integral
    polynomial from both $t_1$ and $z_1$).

    \item Save $z_1$ (currently stored in $w$) into another buffer.

    \item Compute $G$ from the NTRU equation ($G = gF/f \bmod q$) then
    $b = F\adj{f} + G\adj{g}$ (modulo $q$ and modulo $q'$, then
    reassembled as plain 32-bit integers), which is then converted to
    FFT representation and multiplied with $t_1 - z_1$.

    \item Recompute $a = f\adj{f} + g\adj{g}$, first as integers, then
    converted to FFT. This value was already computed previously but we
    did not keep it for RAM-saving reasons. As pointed out previously,
    this recomputation can be done modulo $q'$ if we remember which
    coefficients modulo $q'$ had to be adjusted to yield the plain
    integer value.

    \item Divide the value $(t_1 - z_1)(F\adj{f} + G\adj{g})$ by
    $f\adj{f} + g\adj{g}$. Both polynomials are in FFT representation
    and $f\adj{f} + g\adj{g}$ is self-adjoint, so this uses only
    $n/2$ floating-point inversions. This yields $(t_1 - z_1)L$ (in FFT
    representation).

    \item Set $w$ to $-cF \bmod q$, normalized to coefficients in
    $[-q/2, +q/2]$, and reset $k$ to value $n/2$.

    \item Run $\text{\textsf{ffTreeSampling}}((t_1 - z_1)L, a, W)$.
    We discard the returned value; only $z_0$ (now in buffer $w$)
    will be used afterwards.

    \item Build the second part of the signature $s_1 = f z_0 + F z_1$.
    We do not actually have $z_0$ and $z_1$, since the sampling was
    performed over the fractional parts; we really have obtained
    $z'_0$ and $z'_1$ such that:
    \begin{eqnarray*}
        z_0 &=& \frac{-cF - (-cF \bmod q)}{q} + z'_0 \\
        z_1 &=& \frac{cf - (cf \bmod q)}{q} + z'_1
    \end{eqnarray*}
    This alters the signature computation as follows:
    \begin{eqnarray*}
        s_1 &=& \left(z'_0 + \frac{-cF}{q} - \frac{-cF \bmod q}{q}\right) f
                + \left(z'_1 + \frac{cf}{q} - \frac{cf \bmod q}{q}\right) F \\
            &=& f z'_0 + f \frac{(cF \bmod q)}{q}
                + F z'_1 - F \frac{(cf \bmod q)}{q}
    \end{eqnarray*}
    Note that $s_1$ is integral, hence we can perform that computation
    modulo $q'$; the division by $q$ is then a multiplication by the
    inverse of $q$ modulo $q'$. Working modulo $q'$ is sufficient since
    a valid signature is small ($\norminf{s_1} \leq 840$), hence there
    is no loss of information.

    \item Once $s_1$ is recomputed, the signature verification process is
    performed; this is easier than recomputing $s_0 = c - g z_0 - G z_1$,
    and it more obviously matches what verifiers will compute. The
    verification process needs the public polynomial $h$, which is
    recomputed as $h = g/f \bmod q$ at this point.

    \item If the verification succeeds ($\norm{(s_0, s_1)}$ and
    $\norminf{(s_0, s_1)}$ are low enough) and the signature $s_1$ is
    successfully encoded into bytes within the expected output signature
    size, then a success is reported. Otherwise, the whole process starts
    again ($\mu$ is reused, but a new seed $r$ is generated, hence a new
    polynomial $c$).

\end{enumerate}

The entire procedure can fit in a RAM buffer of size $20n$ bytes (hence
10 or 20 kilobytes, for $n = 512$ or $n = 1024$). In
\textsf{ffTreeSampling} (algorithm~\ref{alg:ffTreeSampling}), the input
uses $12n$ bytes ($8n$ bytes for $t$, $4n$ bytes for $a$, which is
self-adjoint); when doing the first recursive invocation (step 12),
values $t_0$ ($4n$ bytes), $b$ ($2n$ bytes) and $\lambda$ ($2n$ bytes)
must be retained, while for the second recursive invocation (step 15),
values $u_0$ ($4n$ bytes) and $v_1$ ($4n$ bytes) must be retained. In
both cases, the function must reserve $8n$ bytes before calling itself
recursively at half-degree; this leads to a RAM usage of $16n$ bytes by
\textsf{ffTreeSampling}. On top of that, we need two $2n$-byte buffers
for providing the additional input ($-cF \bmod q$, $cf \bmod q$) and
receiving $z_0$ and $z_1$, hence the total RAM usage of $20n$ bytes.

\subsection{Precision}\label{sec:precision}

As pointed out before, both the Falcon tree building and the signature
generation are formally expressed over polynomials in $\bQ[X]/\Phi_n$,
whose exact value can be computed, though this is expensive since the
denominators are huge. Practical implementations of Falcon use
approximations, with floating-point or fixed-point values. The Falcon
security analysis shows that security is maintained as long as the
approximations are ``good enough'', specifically in the inputs to the
\textsf{SamplerZ} algorithm. Each \textsf{SamplerZ} invocation uses as
parameters a centre $x$ and a width $\sigma'$, and returns an integer
sampled along a Gaussian distribution centred on $x$ with standard
deviation $\sigma'$.

In order to estimate how precise our implementation is, we generated 100
signatures with random keys (for both $n = 512$ and $n = 1024$) using
our new implementation. We also generated the same signatures with the
``standard'' implementation, using the same private keys and
pseudorandom seeds, thus resulting in the exact same signature values.
Finally, we computed the mathematically exact rational values for the
inputs to \textsf{SamplerZ} (using a perfunctory Python/Sage
implementation\footnote{At $n = 1024$, this Sage implementation takes
close to 5 minutes to compute one signature, which highlights how
expensive the process is, and why using approximations is necessary.}),
in order to measure how much the two implementations deviate from the
``perfect'' values. At degree $n$, $2n$ centres $x$ and $n$ widths
$\sigma'$ are computed:
\begin{itemize}

    \item For $j = 0$ to $2n - 1$, if the measured centre is $x[j]$, and
    the mathematically exact centre is $x_e[j]$, then the deviation is
    $|x[j] - x_e[j]|$.

    \item For $j = 0$ to $n - 1$, if the measured width is $\sigma'[j]$,
    and the mathematically exact width is $\sigma'_e[j]$, then the deviation
    is $|\sigma'[j]/\sigma'_e[j] - 1|$.

\end{itemize}
Averages of the deviations, for each index $j$, over 100 signatures and
with both the ``standard'' and the new implementations are shown on
figures~\ref{fig:deviations-512} (for $n = 512$) and
\ref{fig:deviations-1024} (for $n = 1024$). Average values are reported
in base-2 logarithmic scale.

\begin{figure}
    \begin{tikzpicture}
        \begin{axis}[
            enlargelimits=false,
            xlabel={Index $j$ (sampling order)},
            ylabel={$\log_2 |x[j] - x_e[j]|$},
            width=0.48\textwidth,
            height=0.48\textwidth,
        ]
            \addplot[
                only marks,
                mark=+,
                color=red,
                mark size=2pt]
            table[meta=dev]{stats_REF512_centre.dat};
            \addplot[
                only marks,
                mark=+,
                color=blue,
                mark size=2pt]
            table[meta=dev]{stats_NEW512_centre.dat};
        \end{axis}
    \end{tikzpicture}
    \hspace{0.04\textwidth}
    \begin{tikzpicture}
        \begin{axis}[
            enlargelimits=false,
            xlabel={Index $j$ (sampling order)},
            ylabel={$\log_2 |\sigma'[j]/\sigma'_e[j] - 1|$},
            width=0.48\textwidth,
            height=0.48\textwidth,
        ]
            \addplot[
                only marks,
                mark=+,
                color=red,
                mark size=2pt]
            table[meta=dev]{stats_REF512_width.dat};
            \addplot[
                only marks,
                mark=+,
                color=blue,
                mark size=2pt]
            table[meta=dev]{stats_NEW512_width.dat};
        \end{axis}
    \end{tikzpicture}
    \caption{\label{fig:deviations-512}Average deviations of centres
    $x$ and widths $\sigma'$ from exact values of the ``standard'' (red)
    and new (blue) implementations ($n = 512$).}
\end{figure}

\begin{figure}
    \begin{tikzpicture}
        \begin{axis}[
            enlargelimits=false,
            xlabel={Index $j$ (sampling order)},
            ylabel={$\log_2 |x[j] - x_e[j]|$},
            width=0.48\textwidth,
            height=0.48\textwidth,
        ]
            \addplot[
                only marks,
                mark=+,
                color=red,
                mark size=2pt]
            table[meta=dev]{stats_REF1024_centre.dat};
            \addplot[
                only marks,
                mark=+,
                color=blue,
                mark size=2pt]
            table[meta=dev]{stats_NEW1024_centre.dat};
        \end{axis}
    \end{tikzpicture}
    \hspace{0.04\textwidth}
    \begin{tikzpicture}
        \begin{axis}[
            enlargelimits=false,
            xlabel={Index $j$ (sampling order)},
            ylabel={$\log_2 |\sigma'[j]/\sigma'_e[j] - 1|$},
            width=0.48\textwidth,
            height=0.48\textwidth,
        ]
            \addplot[
                only marks,
                mark=+,
                color=red,
                mark size=2pt]
            table[meta=dev]{stats_REF1024_width.dat};
            \addplot[
                only marks,
                mark=+,
                color=blue,
                mark size=2pt]
            table[meta=dev]{stats_NEW1024_width.dat};
        \end{axis}
    \end{tikzpicture}
    \caption{\label{fig:deviations-1024}Average deviations of centres
    $x$ and widths $\sigma'$ from exact values of the ``standard'' (red)
    and new (blue) implementations ($n = 1024$).}
\end{figure}

The main observation is that the new implementation consistently
achieves better or at least as good precision as the ``standard'' code.
The precision of the ``standard'' code is already deemed good enough for
the Falcon scheme security goals, with deviations not being exploitable
for attacks within the target security levels; in a sense, it is not
possible to do better than that. However, we can still claim that the
new implementation is also secure, and might have a bit more margin for
alternate ways to compute values.

Why the new implementation seems to have better precision is not
entirely clear, but it is probably due to a combination of these
factors:
\begin{itemize}

    \item Sampling over only the fractional domain implies that the
    centres provided to \textsf{SamplerZ} are lower (in absolute value),
    which, in floating-point formats, means that more bits are available
    for the fractional part.

    \item Computing the input for building the right Falcon sub-tree as
    $d = q^2/(f\adj{f} + g\adj{g})$ entails use of smaller intermediate
    values than the ``standard'' method that uses $F\adj{f} + G\adj{g}$
    and $F\adj{F} + G\adj{G}$, the latter having somewhat large
    coefficients. This, in particular, should explain why the
    ``standard'' implementation has relatively poor precision for the
    sampling widths in the first half of the process (when the right
    sub-tree is used).

    \item The late addition of inputs $(cf \bmod q)/q$ and $(-cF \bmod
    q)/q$ means that these values do not go through a FFT and inverse,
    and thus do not accumulate errors; at the point that they are added
    to the inputs of \textsf{SamplerZ}, they have the maximum possible
    precision that the floating-point format can support.

    \item Generally speaking, floating-point operations return
    approximate values, and errors tend to accumulate; the new
    implementation does fewer floating-point operations, thus
    entailing less error accumulation.

\end{itemize}

\subsection{Reproducibility and Random Hedging}\label{sec:hedging}

As figures~\ref{fig:deviations-512} and \ref{fig:deviations-1024} make
explicit, our new implementation produces centres and widths that are
not exactly the same values as the centres and widths generated by the
``standard'' implementation. In that case, it is unavoidable that even
if using the same pseudorandom seeds for sampling, the two
implementations may diverge: \textsf{SamplerZ} uses rejection sampling,
which necessarily has threshold effects. If a sampled value is below the
threshold for one implementation, but over the threshold for another
implementation for whom the threshold is slightly different, then a
divergence occurs.

To estimate how often divergences occur, we generated one million Falcon
key pairs (at $n = 512$) and computed one signature with each key, with
both the ``standard'' and the new implementations, using identical
seeds. The generated signatures were identical in all cases except 124,
which means that divergences occur at a rate of about $1/8000$, which
interestingly matches the rate observed by \cite{AlmFouLacPre2026}.
Within \textsf{SamplerZ}, the main threshold effect is in the split of
the centre into integral and fractional parts: if the mathematically
exact value is an integer, then each implementation will return an
approximation which may be slightly above or below that integer, the two
cases leading to diverging sampling processes. In the 124 observed
divergences, 121 were in that case, all in either the first two or last
two invocations of \textsf{SamplerZ}. 3 of the divergences, however,
were for centres not particularly near an integer, and came from the
rejecting sampling itself.

The combination of user-provided pseudorandom seeds and divergences
that, while uncommon, still observationally happen, is dangerous: if two
distinct implementations happen to use the same private key over the
same message and with the same seed, and they hit a divergence, then the
observation of the resulting signatures may reveal some information on
the private key. To avoid this issue, our new implementation processes
the user-provided seed in a different way. Namely, for a user-provided
random seed $\rho$, an internal 40-byte seed is derived with value
$H(\text{\textsf{id}} \parallel H(f \parallel g) \parallel \mu \parallel
\rho)$, with $H$ being SHAKE256 (with a 40-byte output), $f$ and $g$
being injected as encoded in the private key, and \textsf{id} being a
conventional fixed string\footnote{This process is in file
\texttt{sign.c}, function \texttt{make\_seedbuf()}.}. In the
``standard'' implementation, \textsf{id} is the empty string; in the new
implementation, it is an 8-byte fixed string. The use of such hashing,
combining the input seed with the private key and the message
representative, is an instance of ``randomness hedging'' and is
considered good security engineering (it maintains security even if the
source of randomness that produced the input seed is of poor quality).
With a different label value in \textsf{id}, it also avoids the
deleterious consequences of two implementations diverging on the same
inputs.

Of course, this different key derivation prevents reproducibility of
test vectors; all the experiments conducted above, to measure precision
and reproducibility, were performed with \textsf{id} set to the same
value (the empty string) as the ``standard'' implementation. In the
published code, the new implementation has its own \textsf{id} and
\emph{none} of the test vectors from the ``standard'' code are
reproduced. As explained above, this is needed for practical security if
deployed applications use both the ``standard'' code and this new
implementation with the same private keys, and inject fixed seeds (which
is certainly not recommended, but may happen as long as the API that
accepts an explicit seed exists).

\subsection{Benchmarks}

Within the course of writing our new implementation, we wrote some
improved functions for computations in $\bZ_q[X]/\Phi_n$, leveraging
SSE2 (on x86) and NEON (on aarch64) SIMD instructions, and using
Montgomery's trick for inverting polynomials in NTT representation with
a single inversion in $\bZ_q$. These optimizations were backported to
the ``standard'' implementation, where they provide some relatively
minor speed-ups.

All our code is constant-time, within the limitations of such exercises
in the presence of modern compilers\cite{Por2025-2}.

\paragraph{RAM usage.}
The ``standard'' code requires a temporary area of size $59n$ bytes;
this is in addition to a stack space of slightly over 1~kB (stack usage
depends on the used compiler). This translates to about 31~kB for $n =
512$, 61~kB for $n = 1024$. The new implementation significantly
improves on that, since it reduces the temporary area size to $20n$
bytes; stack usage has also been slightly reduced (measured at 992 bytes
on our Arm Cortex M4 test platform). In other words, we achieve a RAM
usage of about \textbf{11~kB} (for $n = 512$) or \textbf{21~kB} (for $n
= 1024$). Compared to the ``standard'' implementation, we divided RAM
requirements by a factor of 3. This makes signature generation
substantially more plausible on small embedded systems, including
smartcards, were RAM resources are often very scarce. It may be noted
that this RAM usage is now slightly below that of the key pair
generation (at $22n$ bytes for the temporary area, and less than 1~kB of
stack space).

\paragraph{Speed.}
The test x86 platform is an Intel Core i5-8259U (Coffee Lake, a Skylake
variant) running at 2.3~GHz, in 64-bit mode. TurboBoost is disabled and
the machine (running an up-to-date Ubuntu 24.04) is otherwise idle. The
test aarch64 machine is a Raspberry Pi 5, with a Broadcom BCM2712 CPU
(Arm Cortex-A76, 2.4~GHz); it also runs Ubuntu 24.04. For both machines,
the C compiler is Clang 18.1.3, with optimization flags \verb+-O2+. The
C code on x86 leverages SSE2 but not AVX2 opcodes\footnote{Some
AVX2-specific code is used in key pair generation and in signature
verification, but not in signature generation. Since the x86 ABI only
guarantees presence of SSE2, use of AVX2 requires runtime testing of
support, and code duplication, which is cumbersome from an engineering
perspective. Right now, neither the ``standard'' code nor the new
implementation use AVX2 for signature generation; using AVX2 would
presumably speed things up for both, though it would not make the whole
process twice faster, since at least \textsf{SamplerZ} is
non-parallel.}. To represent embedded systems, we use an STM32F407VGT6
microcontroller (Arm Cortex M4F core); the code is uploaded in the SRAM
block, while the data and stack are located in the CCM area: this
specific setup allows more efficient loading of the code (with no Flash
writes, thus avoiding wearing out the Flash chip) and, more importantly,
running at 168~MHz with no wait states, which speeds up long-running
benchmarks. The Arm Cortex M4F code is mostly in C (compiled with GCC
13.2.1) but with some routines in hand-optimized assembly, in particular
the floating-point operations, as well as the NTT and inverse NTT modulo
$q$ and modulo $q'$. The recursive call in \textsf{ffTreeSampling} is
also written in assembly in order to reduce the amount of stack space
consumed in the recursion.

Table~\ref{tab:speed} shows the achieved speed, expressed in
average clock cycle count for generating a signature, on the three
platforms, with both standard degrees ($n = 512$ and $n = 1024$) and
both implementations.

\begin{table}
    \begin{center}
        {\renewcommand{\arraystretch}{1.1}%
        \setlength{\tabcolsep}{1em}%
        \begin{tabular}{|l|r|r|r|r|}
            \hline
            \textsf{\textbf{Platform}}
            & \multicolumn{2}{c|}{\textsc{Falcon}-512}
            & \multicolumn{2}{c|}{\textsc{Falcon}-1024} \\
            \cline{2-5}
            & \multicolumn{1}{c|}{standard}
            & \multicolumn{1}{c|}{new code}
            & \multicolumn{1}{c|}{standard}
            & \multicolumn{1}{c|}{new code} \\
            \hline
            x86 (64-bit, SSE2)     &   872000 &  1023000 &  1759000 &  2064000 \\
            aarch64 (64-bit, NEON) &   823000 &  1045000 &  1660000 &  2113000 \\
            Arm Cortex M4 (32-bit) & 17020000 & 13450000 & 36570000 & 28580000 \\
            \hline
        \end{tabular}}
    \end{center}
    \caption{\label{tab:speed}Speed benchmarks of the ``standard''
    and new implementations.}
\end{table}

We can observe that the new code is slightly slower than the
``standard'' implementation on large CPUs with efficient hardware
support for floating-point operations (cost is increased by 17\% on x86,
27\% on aarch64); on Arm Cortex M4, though, the reverse happens: the new
code is faster (cost is reduced by 21\%). The new code performs fewer
floating-point operations, but many more integer computations (modulo
$q$ and $q'$), so the overall performance depends on the relative cost
of floating-point and integer operations on the involved hardware. In
our new code, there are no fewer that 22 NTTs (12 modulo $q$, 10 modulo
$q'$) and 12 inverse NTTs (8 modulo $q$, 4 modulo $q'$); we also need to
invert $f$ modulo $q$ twice, and the message representative is derived
into polynomial $c$ three times. These large numbers of operations are
due to the optimization for space, that prevents us from caching
intermediate values across calls to \textsf{ffTreeSampling}. On the
other hand, there are only 3 FFTs, two of which being
\textsf{FFT-selfadj-new} (for $f\adj{f} + g\adj{g}$, which is computed
twice), and no inverse FFT at all.

On Arm Cortex M4, the expensive floating-point operations are additions,
combined addition-subtractions, multiplications, divisions and square
roots, all of which being implemented with hand-optimized assembly
routines (constant-time). A signature generation at $n = 512$ involves
on average 29416 additions, 13322 combined addition-subtractions, 48601
multiplications, 2817 divisions, and 512 square roots; for additions,
addition-subtractions and multiplications, this is less than half the
numbers on the ``standard'' code. In total, these floating-point
operations account for about 6.6 million cycles, i.e. slightly less than
half of the total cost (in the ``standard'' implementation, they make up
to 66\% of the cost). Presumably, a fixed-point implementation would
reduce the cost of these operations, and correspondingly speed up the
signature generation.

\paragraph{Comparison with ML-DSA.}
ML-DSA\cite{Fips204} is the recently NIST-standardized ``generic''
lattice-based post-quantum signature scheme. For signature generation,
the ``\verb+m4f+'' implementation from the pqm4
project\cite{KanRijSchSto2019} computes an ML-DSA-44 signature with an
average cost of 3.943 million cycles on Arm Cortex M4, which is
certainly faster than the 13.45 million cycles of our new Falcon-512
implementation. However, it does that with a much larger RAM space usage
(about 49~kB, compared to 11~kB for our new implementation), which is
likely to hinder deployment on smartcards and similarly small embedded
systems. \cite{BosRenSpr2022} presents a RAM-optimized ML-DSA
implementation which can fit ML-DSA-44 signature generation in 5~kB, but
the runtime cost raises to 18.47 million cycles.

An additional consideration is the variance on signature generation
cost. Both Falcon and ML-DSA may fail to generate a signature at first,
and then need to retry. For Falcon with the current behaviour (and
especially the enforcement of $\norminf{(s_0, s_1)}\leq 840$), this
happens with probability close to $1/2450$ at $n = 512$. This means that
if a signature must be generated within a specific time frame with
probability at least $1 - p$ for some small $p$, then the application
must account for a number of successive attempts. If $p = 10^{-9}$ (for
``nine nines'' reliability) then the application must account for up to
three signing iterations, or about 40 million cycles. For ML-DSA, the
retry probability is much higher: on average, ML-DSA-44 needs 4.25
iterations to succeed, and ``nine nines'' availability needs to account
for no fewer than 78 iterations. The pqm4 benchmarks indicate that along
with the average cost of 3.943 million cycles, the minimal observed cost
is 1.813 million cycles, which would correspond to cases that succeed at
first try. This allows us to estimate that each iteration has cost
$(3.943 - 1.813)/3.25 = 0.655$ million cycles, in which case an ML-DSA
instance that uses 78 iterations would use about 52 million cycles. This
is more than our 40 million cycles, and this is with the ``fast'' ML-DSA
implementation which uses 49~kB of RAM. In that sense, one may argue
that with our implementation, Falcon-512 signature generation on Arm
Cortex M4 is faster \emph{and} smaller than ML-DSA-44. This makes Falcon
quite competitive.

\subsection{Future Work}

The following points need further exploration:
\begin{itemize}

    \item The fixed-point techniques from~\cite{AlmFouLacPre2026}
    should be merged into this new implementation, and benchmarked.
    This should yield some performance improvement on Arm Cortex M4,
    and, more importantly, it would presumably make it easier to
    mask against side-channel attacks.

    \item \textsf{FFT-new} replaces \textsf{FFT-classic}, but the
    \textsf{split} and \textsf{merge} operations from
    \textsf{ffSampling} are unchanged; these operations really are the
    inner primitives of \textsf{FFT-classic} and its inverse. There
    might be possible improvements to this process similar to the
    improvements in \textsf{FFT-new}.

    \item The overall code footprint is non-negligible, at about 53~kB
    for the signature generation code (including 13~kB which are also
    used by key pair generation and verification, in particular the
    SHAKE implementation, and the NTT tables modulo $q$). The
    floating-point code includes 20~kB of constants ($R$ for
    \textsf{FFT-selfadj-new}, $R'$, $C$ and $D$ for \textsf{LDL-new} and
    \textsf{mul-iXadj}, and $\sin(j\pi/1024)$ for $j = 0$ to $511$ for
    \textsf{split} and \textsf{merge}). Embedded systems usually have
    more ROM/Flash than RAM\footnote{1 bit of SRAM needs 6 transistors,
    but 1 bit of ROM or Flash uses about the space of a single
    transistor; moreover, SRAM continuously draws current to maintain
    its contents, proportionally to its size, but ROM and Flash do
    not.}, but this is still a resource that should be preserved, and
    some effort should be made at reducing code size.

\end{itemize}

% ----------------------------------------------------------------------

\newpage
\begin{thebibliography}{20}

% Ideally we would do that only for URLs, but I don't know how to do
% that in a non-painful way.
\RaggedRight

%\bibitem{AbdHwaKanSpr2022}
%A.~Abdulrahman, V.~Hwang, M.~Kannwischer and A.~Sprenkels,
%\emph{Faster Kyber and Dilithium on the Cortex-M4},
%Applied Cryptography and Network Security -- ACNS~2022,
%Lecture Notes In Computer Science, vol.~13269, pp.~853-871, 2022.

%\bibitem{AAPCS}
%\emph{Procedure Call Standard for the Arm® Architecture},
%Arm Limited, version 2024Q3, 2024.

\bibitem{AlmFouLacPre2026}
D.~De Almeida Braga, P.-A.~Fouque, B.~Lachguel and T.~Prest,
\emph{Toward a Secure Fixed-Point Implementation of the Falcon Signature
Scheme},
Advances in Cryptology - CRYPTO 2026 (part III), Lecture Notes in Computer
Science, vol.~16802, pp.~166-197, 2026.

%\bibitem{ChoYooSeo2024}
%J.~Choi, S.~Yoon and S.C.~Seo,
%\emph{Optimized Falcon Verify on Cortex-M4 for Post-Quantum secure UAV 
%communications},
%ICT Express, Elsevier, 2024,\\
%\url{https://doi.org/10.1016/j.icte.2024.11.002}

%\bibitem{DuzFieFis2024}
%S.~Düzlü, R.~Fiedler and M.~Fischlin,
%\emph{BUFFing FALCON without Increasing the Signature Size},\\
%\url{https://eprint.iacr.org/2024/710}

%\bibitem{GajJanKin2024}
%P.~Gajland, J.~Janneck and E.~Kiltz,
%\emph{A Closer Look at Falcon},\\
%\url{https://eprint.iacr.org/2024/1769}

%\bibitem{Hwa2024}
%V.~Hwang,
%\emph{Formal Verification of Emulated Floating-Point Arithmetic in Falcon},\\
%Advances in Information and Computer Security: 19th International Workshop
%on Security -- IWSEC 2024, Lecture Notes in Computer Science,
%vol.~14977, pp.~125-141, 2024.

%\bibitem{ieee754}
%\emph{IEEE Standard for Binary Floating-Point Arithmetic}, 1985,\\
%\url{https://doi.org/10.1109%2FIEEESTD.1985.82928}

%\bibitem{Gol1991}
%D.~Goldberg,
%\emph{What Every Computer Scientist Should Know About Floating-Point
%Arithmetic},
%ACM Computing Surveys (CSUR), vol.~23, issue~1, pp.~5-48, 1991.

%\bibitem{Por2019}
%T.~Pornin,
%\emph{New Efficient, Constant-Time Implementations of Falcon},\\
%\url{https://eprint.iacr.org/2019/893}

%\bibitem{Por2020-1}
%T.~Pornin,
%\emph{Efficient Elliptic Curve Operations On Microcontrollers With
%Finite Field Extensions},\\
%\url{https://eprint.iacr.org/2020/009}

%\bibitem{XKCP}
%G.~Bertoni, J.~Daemen, S.~Hoffert, M.~Peeters, G.~Van~Assche and
%R.~Van~Keer,
%\emph{eXtended Keccak Code Package},\\
%\url{https://github.com/XKCP/XKCP}

%\bibitem{AarAra2022}
%M.~Aardal and D.~Aranha,
%\emph{2DT-GLS: Faster and exception-free scalar multiplication in the
%GLS254 binary curve},\\
%\url{https://eprint.iacr.org/2022/748}

%\bibitem{AmbBosFayJoyLocMur2017}
%C.~Ambrose, J.~Bos, B.~Fay, M.~Joye, M.~Lochter and B.~Murray,
%\emph{Differential Attacks on Deterministic Signatures},\\
%\url{https://eprint.iacr.org/2017/975}

%\bibitem{X962}
%Accredited Standards Committee X9,
%\emph{The Elliptic Curve Digital Signature Algorithm (ECDSA)},
%American National Standards Institute, ANS X9.62, 2005.

%\bibitem{X963}
%Accredited Standards Committee X9,
%\emph{Key Agreement and Key Transport Using Elliptic Curve Cryptography},
%American National Standards Institute, ANS X9.63, 2017.

%\bibitem{AntBroGalLamStrVan2005}
%A.~Antipa, D.~Brown, R.~Gallant, R.~Lambert, R.~Struik and S.~Vanstone,
%\emph{Accelerated Verification of ECDSA signatures},
%Selected Areas in Cryptography - SAC 2005, Lecture Notes in Computer
%Science, vol.~3897, pp.~307-318, 2005.

%\bibitem{AraLopHan2010}
%D.~Aranha, J.~López and D.~Hankerson,
%\emph{Efficient Software Implementation of Binary Field Arithmetic Using
%Vector Instruction Sets},
%Progress in Cryptology - LATINCRYPT 2010,
%Lecture Notes in Computer Science, vol.~6212, pp.~144-161, 2010.

%\bibitem{RistrettoWeb}
%T.~Arcieri, I.~Lovecruft and H.~de Valence,
%\emph{The Ristretto Group},\\
%\url{https://ristretto.group/}

%\bibitem{Atk1992}
%A.~Atkin,
%\emph{Probabilistic primality testing (summary by F. Morain)}, Technical Report 1779, INRIA, 1992,\\
%\url{http://algo.inria.fr/seminars/sem91-92/atkin.pdf}

\bibitem{Bab1986}
L.~Babai,
\emph{On Lovasz’ lattice reduction and the nearest lattice point problem},
Combinatorica, vol.~6, pp.~1-13, 1986.

%\bibitem{BasMaiPop2023}
%A.~Basso, L.~Maino and G.~Pope,
%\emph{FESTA: Fast Encryption from Supersingular Torsion Attacks},\\
%\url{https://ia.cr/2023/660}

%\bibitem{Ber2006}
%D.~Bernstein,
%\emph{Curve25519: new Diffie-Hellman speed records},
%Public Key Cryptography - PKC 2006,
%Lecture Notes in Computer Science, vol.~3958, pp.~207-228, 2006.

%\bibitem{BerDuiLanSchYan2012}
%D.~Bernstein, N.~Duif, T.~Lange, P.~Schwabe and B.-Y.~Yang,
%\emph{High-speed high-security signatures},
%Journal of Cryptographic Engineering, vol.~2, issue~2, pp.~77-89, 2012.

%\bibitem{eBACS}
%D.~Bernstein and T.~Lange,
%\emph{eBACS: ECRYPT Benchmarking of Cryptographic Systems},\\
%\url{https://bench.cr.yp.to} (accessed 4 August 2022).

%\bibitem{BerLanRez2008}
%D.~Bernstein, T.~Lange and R.~Rezaeian Farashahi,
%\emph{Binary Edwards Curves},
%Cryptographic Hardware and Embedded Systems - CHES 2008,
%Lecture Notes in Computer Science, vol.~5154, pp.~244-265, 2008.

%\bibitem{BerHamKraLan2013}
%D.~Bernstein, M.~Hamburg, A.~Krasnova and T.~Lange,
%\emph{Elligator: elliptic-curve points indistinguishable from uniform
%random strings},
%Proceedings of the 2013 ACM SIGSAC Conference on Computer and Communications
%Security, 2013,\\
%\url{https://doi.org/10.1145/2508859.2516734}

%\bibitem{BerYan2019}
%D.~Bernstein and B.-Y.~Yang,
%\emph{Fast constant-time gcd computation and modular inversion},\\
%\url{https://gcd.cr.yp.to/papers.html#safegcd}

\bibitem{Ber2026}
P.-A.~Berthet,
\emph{Masking, Sequences, and FALCON: A Theoretical Study on Masking
Strategies Using Sequences for Non-Linear Operators in the FALCON
Post-Quantum Signature},\\
\url{https://eprint.iacr.org/2026/1534}

\bibitem{BerPaiTavBosCol2025}
P.-A.~Berthet, J.~Paillet, C.~Tavernier, L.~Bossuet and B.~Colombier,
\emph{Masked Computation of the Floor Function and Its Application to
the FALCON Signature},
IACR Communications in Cryptology, vol.~1, issue~4, 2025.

\bibitem{BerTav2025}
P.-A.~Berthet, C.~Tavernier,
\emph{Improving the masked division for the FALCON signature},\\
\url{https://eprint.iacr.org/2025/628}

%\bibitem{BilJoy2003}
%O.~Billet and M.~Joye,
%\emph{The Jacobi model of an elliptic curve and side-channel analysis},
%AAECC-15, Lecture Notes in Computer Science, vol.~2643, pp.~34-42, 2003.

%\bibitem{Bod2007}
%M.~Bodrato,
%\emph{Towards Optimal Toom-Cook Multiplication for Univariate and
%Multivariate Polynomials in Characteristic 2 and 0},
%International Workshop on the Arithmetic of Finite Fields - WAIFI 2007,
%Lecture Notes in Computer Science, vol.~4547, pp.~116-133, 2007.

\bibitem{BosRenSpr2022}
J.~Bos, J.~Renes and A.~Sprenkels,
\emph{Dilithium for Memory Constrained Devices},
Progress in Cryptology - AFRICACRYPT 2022, Lecture Notes in Computer
Science, vol.~13503, pp.~217-235, 2022.

%\bibitem{BreGauThoZim2008}
%R.~Brent, P.~Gaudry, E.~Thomé and P.~Zimmerman,
%\emph{Faster multiplication in GF(2)[x]},
%Algorithmic Number Theory - ANTS-VIII, Lecture Notes in Computer Science,
%vol.~5011, pp.~153-166, 2008.

%\bibitem{BriCorIcaMadRanTib2010}
%É.~Brier, J.-S.~Coron, T.~Icart, D.~Madore, H.~Randriam and M.~Tibouchi,
%\emph{Efficient Indifferentiable Hashing into Ordinary Elliptic Curves},
%Advances in Cryptology - CRYPTO 2010, Lecture Notes in Computer Science,
%vol.~6223, pp.~237-254, 2010.

%\bibitem{Bri1983}
%E.~Brickell,
%\emph{A Fast Modular Multiplication Algorithm with Application to Two Key
%Cryptography},
%Advances in Cryptology - CRYPTO 1982, pp.~51-60, 1983

%\bibitem{BruCurHof1993}
%H.~Brunner, A.~Curriger and M.~Hofstetter,
%\emph{On computing multiplicative inverses in GF($2^m$)},
%IEEE Transactions on Computers, vol.~42, issue~8, pp.~1010-1015, 1993.

%\bibitem{CanKal1991}
%D.~Cantor and E.~Kaltofen,
%\emph{On fast multiplication of polynomials over arbitrary algebras},
%Acta Informatica, vol.~28, issue~7, pp.~693-701, 1991.

%\bibitem{SEC2v2}
%Certicom Research,
%\emph{SEC 2: Recommended Elliptic Curve Domain Parameters},
%version~2.0, 2010,\\
%\url{https://www.secg.org/sec2-v2.pdf}

%\bibitem{Cha2022}
%K.~Chalkias,
%\emph{ed25519-unsafe-libs},\\
%\url{https://github.com/MystenLabs/ed25519-unsafe-libs}

%\bibitem{ChaGarNik2020}
%K.~Chalkias, F.~Garillot and V.~Nikolaenko,
%\emph{Taming the Many EdDSAs},
%Security Standardisation Research - SSR 2020, Lecture Notes in
%Computer Science, vol.~12529, pp.~67-90, 2020.

\bibitem{CheChe2024}
K.-Y.~Chen and J.-P.~Chen,
\emph{Masking Floating-Point Number Multiplication and Addition of Falcon},
IACR Transactions on Cryptographic Hardware and Embedded Systems,
vol.~2024, issue~2, pp.~276-303, 2024.

\bibitem{CheMinLiuGaoZho2025}
S.~Chen, J.~Ming, Y.~Liu, Y.~Gao and Y.~Zhou,
\emph{Masked Gadgets for Integer-Floating-Point Conversion with
Applications to Falcon},
2025 IEEE 43rd International Conference on Computer Design (ICCD),
pp.~507-514, 2025.

%\bibitem{ChuChu1986}
%D.~Chudnovsky and G.~Chudnovsky,
%\emph{Sequences of numbers generated by addition
%in formal groups and new primality and factorization tests},
%Advances in Applied Mathematics, vol.~7, issue~4, pp.~385–434, 1986.

%\bibitem{CohFre2006}
%H.~Cohen and G.~Frey (editors),
%\emph{Handbook of Elliptic and Hyperelliptic Curve Cryptography},
%Chapman and Hall/CRC, 2006.

%\bibitem{Col1726}
%J.~Colson,
%\emph{A Short Account of Negativo-Affirmative Arithmetick},
%Philosophical Transactions (1683-1775), vol.~34, pp.~161-173, 1726.

%\bibitem{FrostDraft}
%D.~Connolly, C.~Komlo, I.~Goldberg and C.~Wood,
%\emph{Two-Round Threshold Schnorr Signatures with FROST},\\
%\url{https://datatracker.ietf.org/doc/html/draft-irtf-cfrg-frost-08}

\bibitem{CooTuk1965}
J.~Cooley and J.~Tukey,
\emph{An algorithm for the machine calculation of complex Fourier series},
Mathematics of Computation, vol.~19, issue~90, pp.~297-301.

%\bibitem{CreJac2019}
%C.~Cremers and D.~Jackson,
%\emph{Prime, Order Please! Revisiting Small Subgroup and Invalid Curve
%Attacks on Protocols using Diffie-Hellman},
%IEEE 32nd Computer Security Foundations Symposium (CSF), 2019.

%\bibitem{dalek-crypto}
%I.~Lovecruft and H.~de Valence,
%\emph{Dalek cryptography},\\
%\url{https://github.com/dalek-cryptography}

%\bibitem{DifHel1976}
%W.~Diffie and M.~Hellman,
%\emph{New directions in cryptography},
%IEEE Transactions on Information Theory, vol.~22, issue~6, pp.~644-654, 1976.

\bibitem{DucPre2016}
L.~Ducas and T.~Prest,
\emph{Fast Fourier Orthogonalization},
ISSAC '16: Proceedings of the 2016 ACM International Symposium on Symbolic
and Algebraic Computation, pp.~191-198, 2016.

%\bibitem{Duq2007}
%S.~Duquesne,
%\emph{Improving the arithmetic of elliptic curves in the Jacobi model},
%Information Processing Letters, vol.~104, issue~3, pp.~101–105, 2007.

%\bibitem{FisPat1974}
%M.~Fischer and M.~Paterson,
%\emph{String-Matching and Other Products},
%MAC Technical Memorandum 41, Massachusetts Institute of Technology, 1974.

%\bibitem{Fog2023}
%A.~Fog,
%\emph{The microarchitecture of Intel, AMD, and VIA CPUs},\\
%\url{https://www.agner.org/optimize/microarchitecture.pdf}

\bibitem{FouGajGroJanKil2026}
P.-A.~Fouque, P.~Gajland, H.~de~Groote, J.~Janneck and E.~Kiltz,
\emph{A Closer Look at Falcon},
Advances in Cryptology - EUROCRYPT 2026 (part IV), Lecture Notes in
Computer Science, vol.~16544, pp.~3-31, 2026.

%\bibitem{GalLamVan2001}
%R.~Gallant, J.~Lambert and S.~Vanstone,
%\emph{Faster Point Multiplication on Elliptic Curves with Efficient
%Endomorphisms},
%Advances in Cryptology - CRYPTO 2001, Lecture Notes in Computer Science,
%vol.~2139, pp.~190-200, 2001.

%\bibitem{GalLinSco2009}
%S.~Galbraith, X.~Lin and M.~Scott,
%\emph{Endomorphisms for Faster Elliptic Curve Cryptography on a Large
%Class of Curves},
%Advances in Cryptology - EUROCRYPT 2009, Lecture Notes in Computer Science,
%vol.~5479, pp.~518-535, 2009.

%\bibitem{GauHesSma2002}
%P.~Gaudry, F.~Hess and N.~Smart,
%\emph{Constructive and destructive facets of Weil descent on elliptic curves},
%Journal of Cryptology, vol.~15, issue~1, pp.~19-46, 2002.

\bibitem{Gau1866}
C.~Gauss,
\emph{Nachlass: Theoria interpolationis methodo nova tractata},
Werke, vol.~3, pp.~265-303, Göttingen: Königliche Gesellschaft der Wissenschaften, 1866.

\bibitem{GenSan1966}
W.~Gentleman and G.~Sande,
\emph{Fast Fourier transforms -- for fun and profit},
AFIPS '66 (Fall), pp.~563-578, 1966.

%\bibitem{Ham2009}
%M.~Hamburg,
%\emph{Accelerating AES with Vector Permute Instructions},
%Cryptographic Hardware and Embedded Systems - CHES 2009,
%Lecture Notes in Computer Science, vol.~5747, pp.~18-32, 2009.

%\bibitem{Ham2015-1}
%M.~Hamburg,
%\emph{Decaf: Eliminating cofactors through point compression},
%Advances in Cryptology - CRYPTO 2015, Lecture Notes in Computer Science,
%vol.~9215, pp.~705-723, 2015.

%\bibitem{Ham2015-2}
%M.~Hamburg,
%\emph{Ed448-Goldilocks, a new elliptic curve},\\
%\url{https://eprint.iacr.org/2015/625}

%\bibitem{HanKarMen2009}
%D.~Hankerson, K.~Karabina and A.~Menezes,
%\emph{Analyzing the Galbraith-Lin-Scott Point Multiplication Method for
%Elliptic Curves over Binary Fields},
%IEEE Transactions on Computers, vol.~58, issue~10, pp.~1411-1420, 2009.

%\bibitem{HanMenVan2004}
%D.~Hankerson, A.~Menezes and S.~Vanstone,
%\emph{Guide to Elliptic Curve Cryptography},
%Springer/Sci-Tech/Trade, 2004.

%\bibitem{HisWonCarDaw2008}
%H.~Hisil, K.~Wong, G.~Carter and E.~Dawson,
%\emph{Twisted Edwards Curves Revisited},
%Advances in Cryptology - ASIACRYPT 2008, Lecture Notes in Computer Science,
%vol.~5350, pp.~326-343, 2008.

%\bibitem{HisWonCarDaw2009}
%H.~Hisil, K.~Wong, G.~Carter and E.~Dawson,
%\emph{Jacobi Quartic Curves Revisited},
%Information Security and Privacy - ACISP 2009, Lecture Notes in Computer
%Science, vol.~5594, pp.~452-468, 2009.

%\bibitem{Hon1989}
%J.~Honkala,
%\emph{On Number Systems with Negative Digits},
%Annales Academiæ Scientiarum Fennicæ, Series A. I. Mathematica,
%vol.~14, issue~1, pp.149-156, 1989.

%\bibitem{ItoTsu1988}
%T.~Itoh and S.~Tsujii,
%\emph{A Fast Algorithm for Computing Multiplicative Inverses in
%$\GF(2^m)$ Using Normal Bases},
%Information and Computation, vol.~78, pp.~171-177, 1988.

%\bibitem{Jac1829}
%C.~G.~J.~Jacobi,
%\emph{Fundamenta nova theoriae functionum ellipticarum},
%Sumtibus fratrum, 1829.

%\bibitem{EdDSArfc8032}
%S.~Josefsson and I.~Liusvaara,
%\emph{Edwards-Curve Digital Signature Algorithm (EdDSA)},\\
%\url{https://tools.ietf.org/html/rfc8032}

\bibitem{KanRijSchSto2019}
M.~Kannwischer, J.~Rijneveld, P.~Schwabe and K.~Stoffelen,
\emph{pqm4: Testing and Benchmarking NIST PQC on ARM Cortex-M4},\\
\url{https://eprint.iacr.org/2019/844}

\bibitem{KarAys2024}
E.~Karabulut and A.~Aysu,
\emph{Masking FALCON’s floating-point multiplication in hardware},
IACR Transactions on Cryptographic Hardware and Embedded Systems,
vol.~2024, issue~4, pp.~483-508, 2024.

%\bibitem{KarOfm1962}
%А.~Карацуба and Ю.~Офман,
%\emph{Умножение Многозначных Чисел На Автоматах},
%Доклады Академи и наук СССР, vol.~145, issue~2, pp.~293-294, 1962.
%{\fontencoding{T2A}\fontfamily{NimbusSerif}\selectfont А.~Карацуба}
%and {\fontencoding{T2A}\fontfamily{NimbusSerif}\selectfont Ю.~Офман,
%\emph{Умножение Многозначных Чисел На Автоматах},
%Доклады Академи и наук СССР}, vol.~145, pp.~293-294, 1962.
%A.~Karatsuba and Y.~Ofman,
%\emph{Multiplication of Many-Digital Numbers by Automatic Computers},
%Proceedings of the USSR Academy of Sciences, vol.~145, pp.~293-294, 1962.

%\bibitem{KimKim2007}
%K.~Kim and S.~Kim,
%\emph{A New Method for Speeding Up Arithmetic on Elliptic Curves over
%Binary Fields},\\
%\url{https://eprint.iacr.org/2007/181}

\bibitem{Kle2000}
P.~Klein,
\emph{Finding the closest lattice vector when it’s unusually close},
SODA '00: Proceedings of the eleventh annual ACM-SIAM symposium on
Discrete algorithms, pp.~937-941, 2000.

%\bibitem{Knu1969}
%D.~Knuth,
%\emph{The Art of Computer Programming, volume~2: Seminumerical Algorithms},
%Addison-Wesley, 1969.

%\bibitem{KomGol2020}
%C.~Komlo and I.~Goldberg,
%\emph{FROST: Flexible Round-Optimized Schnorr Threshold Signatures},\\
%\url{https://eprint.iacr.org/2020/852}

%\bibitem{rfc7748}
%A.~Langley, M.~Hamburg and S.~Turner,
%\emph{Elliptic Curves for Security},\\
%\url{https://www.rfc-editor.org/rfc/rfc7748}

%\bibitem{Lean}
%{\fontspec{texgyreadventor-regular.otf}\selectfont{L\kern -0.1em\reflectbox{E}\kern -0.13em\raisebox{\depth}{\scalebox{1}[-1]{A}}\kern -0.13em{N}}} \emph{Programming Language},\\
%\url{https://lean-lang.org/}

%\bibitem{Len2019}
%E.~Lenngren,
%\emph{X25519 for ARM AArch64},\\
%\url{https://github.com/Emill/X25519-AArch64}

%\bibitem{Len2020}
%E.~Lenngren,
%\emph{X25519 for ARM Cortex-M4 and other ARM processors},\\
%\url{https://github.com/Emill/X25519-Cortex-M4}

%\bibitem{LopDah1999}
%J.~López and R.~Dahab,
%\emph{Fast Multiplication on Elliptic Curve Over $\GF(2^m)$ without
%precomputation},
%Cryptographic Hardware and Embedded Systems - CHES'99,
%Lecture Notes in Computer Science, vol.~1717, pp.~316-327, 1999.

%\bibitem{MenOorVan1996}
%A.~Menezes, P.~van Oorschot and S.~Vanstone,
%\emph{Handbook of Applied Cryptography},
%CRC Press, 1996,\\
%\url{https://cacr.uwaterloo.ca/hac/}

%\bibitem{Mon1985}
%P.~Montgomery,
%\emph{Modular Multiplication Without Trial Division},
%Mathematics of Computation, vol.~44, issue~170, pp.~519-521, 1985.

%\bibitem{MooShu2011}
%D.~Moody and D.~Shumow,
%\emph{Analogues of Vélu's Formulas for Isogenies on Alternate Models of
%Elliptic Curves},
%\url{https://eprint.iacr.org/2011/430}

%\bibitem{NacSte2001}
%D.~Naccache and J.~Stern,
%\emph{Signing on a Postcard},
%Financial Cryptography - FC 2000, Lecture Notes in Computer Science,
%vol.~1962, pp.~121-135, 2001.

%\bibitem{NatSar2020}
%K.~Nath and P.~Sarkar,
%\emph{Efficient 4-way Vectorizations of the Montgomery Ladder},\\
%\url{https://eprint.iacr.org/2020/378}

%\bibitem{NevSmaWar2009}
%G.~Neven, N.~P.~Smart and B.Warinschi,
%\emph{Hash function requirements for Schnorr signatures},
%Journal of Mathematical Cryptology, vol.~3, issue 1, pp.~69-87, 2009.

\bibitem{NISTPQC}
\emph{Post-Quantum Cryptography},
National Institute of Standard and Technology,\\
\url{https://www.nist.gov/pqcrypto}

%\bibitem{Fips180}
%Information Technology Laboratory,
%\emph{Secure Hash Standard (SHS)},
%National Institute of Standard and Technology, FIPS~180-4, 2015.

%\bibitem{Fips186}
%Information Technology Laboratory,
%\emph{Digital Signature Standard (DSS)},
%National Institute of Standard and Technology, FIPS~186-5, 2023.

%\bibitem{Fips202}
%Information Technology Laboratory,
%\emph{SHA-3 Standard: Permutation-Based Hash and Extendable-Output
%Functions},
%National Institute of Standard and Technology, FIPS~202, 2015.

\bibitem{Fips203}
Information Technology Laboratory,
\emph{Module-Lattice-Based Key-Encapsulation Mechanism Standard},
National Institute of Standard and Technology, FIPS~203, 2024.

\bibitem{Fips204}
Information Technology Laboratory,
\emph{Module-Lattice-Based Digital Signature Standard},
National Institute of Standard and Technology, FIPS~204, 2024.

%\bibitem{BLAKE3}
%J.~O'Connor, J.-P.~Aumasson, S.~Neves and Z.~Wilcox-O'Hearn,
%\emph{BLAKE3},\\
%\url{https://github.com/BLAKE3-team/BLAKE3}

%\bibitem{OliLopAraRod2014}
%T.~Oliveira, J.~López-Hernández, D.~Aranha and F.~Rodríguez-Henríquez,
%\emph{Two is the fastest prime: lambda coordinates for binary elliptic
%curves},
%Journal of Cryptographic Engineering, vol.~4, issue~1, pp.~3-17, 2014.

%\bibitem{OliLopAraRod2016}
%T.~Oliveira, J.~López-Hernández, D.~Aranha and F.~Rodríguez-Henríquez,
%\emph{Improving the performance of the GLS254},
%Presented at the CHES 2016 Rump Session.

%\bibitem{PetQui2012}
%C.~Petit and J.-J.~Quisquater,
%\emph{On Polynomial Systems Arising from a Weil Descent},
%Advances in Cryptology - ASIACRYPT 2012,
%Lecture Notes in Computer Science, vol.~7658, pp.~451-466, 2012.

%\bibitem{PhaKor1994}
%D.~Phatak and I.~Koren,
%\emph{Hybrid Signed-Digit Number Systems: A Unified Framework for
%Redundant Number Representations With Bounded Carry Propagation Chains},
%IEEE Transactions on Computers, vol.~43, issue~8, pp.~880-891, 1994.

%\bibitem{PodSomSchLocRos2018}
%D.~Poddebniak, J.~Somorovsky, S.~Schinzel, M.~Lochter and P.~Rösler,
%\emph{Attacking Deterministic Signature Schemes Using Fault Attacks},
%2018 IEEE European Symposium on Security and Privacy (EuroS\&P),
%pp.~338-352, 2018.

%\bibitem{DeterministicECDSArfc}
%T.~Pornin,
%\emph{Deterministic Usage of the Digital Signature Algorithm (DSA) and
%Elliptic Curve Digital Signature Algorithm (ECDSA)},\\
%\url{https://tools.ietf.org/html/rfc6979}

%\bibitem{BearSSL}
%T.~Pornin,
%\emph{BearSSL: a smaller SSL/TLS library},\\
%\url{https://www.bearssl.org/}

%\bibitem{Por2020-2}
%T.~Pornin,
%\emph{Optimized Lattice Basis Reduction In Dimension 2, and Fast Schnorr
%and EdDSA Signature Verification},\\
%\url{https://eprint.iacr.org/2020/454}

%\bibitem{Por2020-3}
%T.~Pornin,
%\emph{Optimized Binary GCD for Modular Inversion},\\
%\url{https://eprint.iacr.org/2020/972}

%\bibitem{Por2020-5}
%T.~Pornin,
%\emph{Double-Odd Elliptic Curves},\\
%\url{https://eprint.iacr.org/2020/1558}

%\bibitem{Por2022-4}
%T.~Pornin,
%\emph{Double-Odd Jacobi Quartic},\\
%\url{https://eprint.iacr.org/2022/1052}

%\bibitem{Por2022-6}
%T.~Pornin,
%\emph{Efficient and Complete Formulas for Binary Curves},\\
%\url{https://eprint.iacr.org/2022/1325}

%\bibitem{Por2023-3}
%T.~Pornin,
%\emph{On Multiplications with Unsaturated Limbs},\\
%\url{https://research.nccgroup.com/2023/09/18/on-multiplications-with-unsaturated-limbs/}

\bibitem{PorPre2019}
T.~Pornin and T.~Prest,
\emph{More Efficient Algorithms for the NTRU Key Generation using the
Field Norm},
Public Key Cryptography - PKC 2019 (part II), Lecture Notes in Computer
Science, vol.~11443, pp.~504-533, 2019.

\bibitem{Por2023-1}
T.~Pornin,
\emph{Improved Key Pair Generation for Falcon, BAT and Hawk},\\
\url{https://eprint.iacr.org/2023/290}

\bibitem{Por2025-2}
T.~Pornin,
\emph{Constant-Time Code: The Pessimist Case},\\
\url{https://eprint.iacr.org/2025/435}

\bibitem{Por2025-4}
T.~Pornin,
\emph{Improved (Again) Key Pair Generation for Falcon, BAT and Hawk},\\
\url{https://eprint.iacr.org/2025/1239}

\bibitem{FalconRound3}
T.~Prest, P.-A.~Fouque, J.~Hoffstein, P.~Kirchner,
V.~Lyubashevsky, T.~Pornin, T.~Ricosset, G.~Seiler, W.~Whyte and Z.~Zhang,
\emph{Falcon},
Technical report, National Institute of Standards and Technology, 2019,
\url{https://csrc.nist.gov/Projects/post-quantum-cryptography/post-quantum-cryptography-standardization/round-3-submissions}

%\bibitem{RenCosBat2016}
%J.~Renes, C.~Costello and L.~Batina,
%\emph{Complete Addition Formulas for Prime Order Elliptic Curves},
%Advances in Cryptology - EUROCRYPT 2016, Lecture Notes in Computer Science,
%vol.~9665, pp.~403-428, 2016.\\
%\url{https://ia.cr/2015/1060}

%\bibitem{RivShaAdl1978}
%R.~Rivest, A.~Shamir and L.~Adleman,
%\emph{A Method for Obtaining Digital Signatures and Public-Key Cryptosystems},
%Communications of the ACM, vol.~21, issue~2, pp.~120-126, 1978.

%\bibitem{BLAKE2}
%M.-J.~Saarinen and J.-P.~Aumasson,
%\emph{The BLAKE2 Cryptographic Hash and Message Authentication Code (MAC)}\\
%\url{https://datatracker.ietf.org/doc/html/rfc7693}

%\bibitem{Sch1989}
%C.~Schnorr,
%\emph{Efficient identification and signatures for smart cards},
%Advances in Cryptology - CRYPTO '89, Lecture Notes in Computer Science,
%vol.~435, pp.~239-252, 1990.

%\bibitem{ShaWoe2006}
%A.~Shallue and C.~van de Woestijne,
%\emph{Construction of rational points on elliptic curves over finite
%fields},
%Algorithm Number Theory Symposium - ANTS 2006, Lecture Notes in Computer
%Science, vol.~4076, pp.~510-524, 2006.

%\bibitem{Sil2009}
%J.~Silverman,
%\emph{The Arithmetic of Elliptic Curves} (2nd edition),
%Springer New York, NY, 2009.

%\bibitem{Sol2000}
%J.~Solinas,
%\emph{Efficient Arithmetic on Koblitz Curves},
%Designs, Codes and Cryptography, vol.~19, issue~2-3, pp.~195-249, 2000.

%\bibitem{Str1964}
%E.~Straus,
%\emph{Addition chains of vectors (problem 5125)},
%American Mathematical Monthly, vol.~70, pp.~806-808, 1964.

\bibitem{SunZhoZhaTaoQiaMin2022}
S.~Sun, Y.~Zhou, R.~Zhang, Y.~Tao, Z.~Qiao and J.~Ming,
\emph{Fast Fourier orthogonalization over NTRU lattices},
Information and Communications Security - ICICS 2022,
Lecture Notes in Computer Science, vol.~13407, pp.~109-127, 2022.

%\bibitem{Too1963}
%А.~Тоом,
%\emph{О сложности схемы из функциональных злементов, реализирующей
%умножение целых чисел},
%Доклады Академи и наук СССР, vol.~153, issue~3, pp.~496-498, 1963.

%\bibitem{Ula2007}
%M.~Ulas,
%\emph{Rational Points on Certain Hyperelliptic Curves over Finite Fields},
%Bulletin of the Polish Academy of Sciences - Mathematics, vol.~55,
%issue~2, pp.~97-104, 2007.

%\bibitem{RistrettoDraft}
%H.~de Valence, J.~Grigg, M.~Hamburg, I.~Lovecruft, G.~Tankersly and
%F.~Valsorda,
%\emph{The ristretto255 and decaf448 Groups},\\
%\url{https://datatracker.ietf.org/doc/html/draft-irtf-cfrg-ristretto255-decaf448-03}

%\bibitem{Vel1971}
%J.~Vélu,
%\emph{Isogénies entre courbes elliptiques},
%C.R. Acad. Sc. Paris, Série A, vol.~273, pp.~238-241, 1971.

%\bibitem{Was2008}
%L.~Washington,
%\emph{Elliptic Curves: Number Theory and Cryptography} (2nd edition),
%Chapman and Hall/CRC, 2008.

%\bibitem{WhiWat1927}
%E.~Whittaker and G.~Watson,
%\emph{A Course of Modern Analysis},
%Cambridge University Press, 1927.

%\bibitem{Wycheproof}
%\emph{Project Wycheproof},\\
%\url{https://github.com/C2SP/wycheproof}

\end{thebibliography}

% ----------------------------------------------------------------------

\end{document}
